\chapter{Rough Paths} \label{ch:roughpaths}

\section{Introduction}\label{sec:roughintro}

Remark \ref{rem:criticalexponent} closed Part \ref{part:part1} with a diagnosis rather than a
theorem. Brownian motion has finite $p$-variation exactly for $p>2$; Young's integral requires
$p<2$; and the Itô integral of Section \ref{sec:stochasticcalculus} bridges the gap not by refining
an estimate but by \emph{spending probabilistic structure} --- adaptedness, an $L^2(\Omega)$ limit
in place of a pathwise one, and a definition valid only up to null sets. That bargain is what
Part \ref{part:part1} was built on, and it is worth being precise about the two circumstances in
which it collapses.

The first is that the driving signal may be rougher than Brownian. Itô's construction is not a
general theory of integration against irregular paths: it is a theory of integration against
\emph{semimartingales}, in the sense of Definition \ref{def:semimartingale}. The moment the
integrator leaves that class, the entire apparatus --- the isometry of Theorem \ref{thm:itoisometry}, the martingale
property, the localisation of Example \ref{example_stoppingTimes} --- becomes unavailable at once,
and not merely inconvenient.

The second is subtler and, in the end, more consequential. Even where Itô's theory applies, it
delivers the solution of a stochastic differential equation as a measurable function on $\Om$,
constructed through an $L^2$ limit that depends on the filtration and the measure. It does not
tell us that the solution depends \emph{continuously} on the driving trajectory, and in fact, as
we shall see in Subsection \ref{subsec:whatfails}, no such statement is true in the topology one
would naïvely choose. This is not a technical blemish. A model whose output is a discontinuous
function of its input is a model one cannot approximate, cannot discretise with confidence, and
cannot calibrate stably.

Rough path theory, introduced by Lyons \cite{lyons1998}, resolves both difficulties simultaneously,
and by a single device: it enlarges the notion of ``path''. The remainder of this chapter develops
that idea. We begin, however, with the process that forces the question upon us.

\subsection{Fractional Brownian Motion}\label{subsec:fbm}

Brownian motion is characterised, among Gaussian processes with stationary increments, by the
independence of those increments. Relaxing exactly that requirement --- and nothing else --- yields
the following one-parameter family, introduced by Kolmogorov and given its modern form by Mandelbrot
and Van Ness \cite{mandelbrotvanness1968}.

\begin{definition}[Fractional Brownian motion]\label{def:fbm}
Let $\Hurst \in (0,1)$. A \textbf{fractional Brownian motion} with Hurst parameter $\Hurst$ is a
centred Gaussian process $\BH = \{\BH_t\}_{t\geq0}$ with $\BH_0 = 0$ almost surely and covariance
\begin{equation}\label{eq:fbmcovariance}
\E\bigl[\BH_s \BH_t\bigr] \;=\; R_\Hurst(s,t) \;:=\; \tfrac12\Bigl( s^{2\Hurst} + t^{2\Hurst} - |t-s|^{2\Hurst} \Bigr),
\qquad s,t \geq 0 .
\end{equation}
\end{definition}

That \eqref{eq:fbmcovariance} is a legitimate covariance --- that is, symmetric and positive
semi-definite --- is a genuine constraint on $\Hurst$, and it is the reason the parameter range is
$(0,1)$ and not larger; a proof is in \cite[Ch.\ 7]{marcusrosen2006}. Granting it, the process
exists by the standard Gaussian construction, and Kolmogorov's criterion (Theorem
\ref{thm:kolmogorovcontinuity}) supplies a continuous modification, with which we always work in
accordance with Remark \ref{rem:modification}.

\begin{proposition}[Elementary properties]\label{prop:fbmproperties}
Let $\BH$ be a fractional Brownian motion with Hurst parameter $\Hurst \in (0,1)$. Then:
\begin{itemize}
    \item[a)] $\BH$ has \emph{stationary increments}, with
    \begin{equation}\label{eq:fbmincrement}
        \E\bigl[\,|\BH_t - \BH_s|^2\,\bigr] = |t-s|^{2\Hurst}, \qquad \text{so } \BH_t - \BH_s \sim \Normal\bigl(0, |t-s|^{2\Hurst}\bigr).
    \end{equation}
    \item[b)] $\BH$ is \emph{self-similar of index $\Hurst$}: for every $\lambda>0$, the processes
    $\{\BH_{\lambda t}\}_{t\geq0}$ and $\{\lambda^{\Hurst}\BH_t\}_{t \geq 0}$ have the same law.
    \item[c)] Almost every sample path is locally $\gamma$-Hölder continuous for every
    $\gamma < \Hurst$, and for no $\gamma \geq \Hurst$.
    \item[d)] Almost surely, $V_p(\BH)_t < \infty$ for every $p > 1/\Hurst$ and $V_p(\BH)_t = \infty$
    for every $p \leq 1/\Hurst$: the critical exponent of Remark \ref{rem:criticalexponent} is
    $1/\Hurst$.
    \item[e)] $\Hurst = \tfrac12$ recovers Brownian motion. For $\Hurst > \tfrac12$ the increments
    over disjoint intervals are positively correlated, for $\Hurst < \tfrac12$ negatively, and in
    neither case are they independent.
\end{itemize}
\end{proposition}

\textbf{\textit{Proof:}} a) Expanding and using \eqref{eq:fbmcovariance},
\[
\E\bigl[(\BH_t-\BH_s)^2\bigr] = R_\Hurst(t,t) - 2R_\Hurst(s,t) + R_\Hurst(s,s)
= t^{2\Hurst} - \bigl(s^{2\Hurst}+t^{2\Hurst}-|t-s|^{2\Hurst}\bigr) + s^{2\Hurst} = |t-s|^{2\Hurst}.
\]
Since the process is centred and Gaussian, the law of an increment is determined by this variance
alone, and depends on $s,t$ only through $t-s$.

b) $R_\Hurst(\lambda s, \lambda t) = \lambda^{2\Hurst}R_\Hurst(s,t)$ is immediate from
\eqref{eq:fbmcovariance}. Two centred Gaussian processes with the same covariance have the same law.

c) The Gaussian moments give $\E[|\BH_t-\BH_s|^{2m}] = c_m|t-s|^{2m\Hurst}$, so
\eqref{eq:kolmogorovmoment} holds with $\alpha = 2m$ and $\beta = 2m\Hurst - 1$, and Theorem
\ref{thm:kolmogorovcontinuity} yields Hölder exponents $\gamma < \Hurst - 1/2m$, which approach
$\Hurst$ as $m\to\infty$. The negative half is the exact analogue of Theorem \ref{thm:bmregularity}
i) and is proved by the same Paley--Wiener--Zygmund argument, adapted in \cite[Ch.\ 7]{marcusrosen2006}.

d) Follows from c) together with the same competition between $p$-variation and modulus of
continuity that produced Proposition \ref{prop:bmunboundedvariation}; see \cite[Ch.\ 15]{frizvictoir2010}
for the details, which we do not repeat.

e) For $\Hurst=\tfrac12$, \eqref{eq:fbmcovariance} collapses to $\min\{s,t\}$, which is Proposition
\ref{prop:bm-properties} a); by Definition \ref{def:Brownian-motion} and the Gaussian
characterisation this identifies the process. For the sign of the correlation, take
$0\leq s\leq t\leq u\leq v$ and compute directly from \eqref{eq:fbmcovariance},
\[
2\,\E\bigl[(\BH_t-\BH_s)(\BH_v-\BH_u)\bigr] = |v-s|^{2\Hurst} + |u-t|^{2\Hurst} - |u-s|^{2\Hurst} - |v-t|^{2\Hurst},
\]
which is the second difference of the function $\tau \mapsto \tau^{2\Hurst}$, hence positive when
that function is convex ($\Hurst>\tfrac12$) and negative when it is concave ($\Hurst<\tfrac12$).
It vanishes identically only at $\Hurst = \tfrac12$. \qedfilledtext

\medskip

Part e) is the modelling content of the definition and deserves its name. For $\Hurst > \tfrac12$
the process exhibits \emph{persistence}: an increase tends to be followed by an increase, and the
correlations decay so slowly that they are not summable --- the phenomenon known as \emph{long
memory}. For $\Hurst < \tfrac12$ it exhibits \emph{anti-persistence}, or mean reversion at every
scale, and the paths are correspondingly rougher than Brownian. Part d) makes ``rougher'' precise:
the critical exponent $1/\Hurst$ exceeds $2$ exactly when $\Hurst < \tfrac12$.

That last observation is not a curiosity but an obstruction, and we state it as such.

\begin{theorem}[Rogers]\label{thm:fbmnotsemimartingale}
Fractional Brownian motion is a semimartingale in the sense of Definition \ref{def:semimartingale}
if and only if $\Hurst = \tfrac12$. \qedopentext
\end{theorem}

The proof is in \cite{rogers1997}, and the mechanism is worth recording even though we do not
reproduce the argument. For $\Hurst > \tfrac12$ the quadratic variation of Definition
\ref{def:quadratic-variation} vanishes, while the total variation is infinite by part d); a
continuous semimartingale with zero quadratic variation must be of finite variation, by the
uniqueness of the decomposition, and this is a contradiction. For $\Hurst < \tfrac12$ the quadratic
variation is instead infinite, which is incompatible with the decomposition from the other
direction. Only the balance point $\Hurst = \tfrac12$ escapes.

The consequence is severe and complete. \textbf{For $\Hurst \neq \tfrac12$ there is no Itô integral
against $\BH$}, no Itô formula, no martingale representation, and no Girsanov theorem --- not
because the constructions are difficult but because every one of them takes the semimartingale
property as its hypothesis. Definition \ref{def:semimartingale} therefore marks the outer limit of
the theory developed in Chapter \ref{ch:foundations}, and we have now crossed it.

\subsection{Volatility is Rough}\label{subsec:roughvol}

Why should one wish to cross it? The answer, for this dissertation, is empirical, and although a
thesis in mathematics is not the place for an econometric study, the finding that motivates
everything below can be stated in one paragraph.

Let $\sigma_t$ denote spot volatility, estimated from high-frequency returns over a grid of
spacing $\Delta$, and consider the $q$-th absolute moment of the increments of its logarithm,
\begin{equation}\label{eq:scalingmoment}
m(q,\Delta) \;:=\; \frac{1}{N}\sum_{i=1}^{N} \bigl|\log\sigma_{i\Delta} - \log\sigma_{(i-1)\Delta}\bigr|^{q} .
\end{equation}
Gatheral, Jaisson and Rosenbaum \cite{gatheraljaissonrosenbaum2018} observe that, across equity
indices and thousands of individual assets, $m(q,\Delta) \propto \Delta^{\zeta_q}$ with
$\zeta_q = q\Hurst$ --- a linear, \emph{monofractal} scaling with no curvature in $q$ --- and that
the increments are close to Gaussian. Comparison with \eqref{eq:fbmincrement} is immediate: these
are exactly the scaling relations of a fractional Brownian motion. The estimated exponent is
$\Hurst \approx 0.1$, and it is remarkably stable across assets and across decades.

\figurehere{Log--log plot of $m(q,\Delta)$ against $\Delta$ for several values of $q$, showing the
parallel straight lines whose slopes are $\zeta_q$, alongside the plot of $\zeta_q$ against $q$
exhibiting the linear (monofractal) relation $\zeta_q = q\Hurst$. This is the visual statement of
the empirical finding of \cite{gatheraljaissonrosenbaum2018}.}

This inverts an earlier consensus. The fractional stochastic volatility model of Comte and Renault
\cite{comterenault1998} also used fractional Brownian motion, but with $\Hurst > \tfrac12$, in order
to reproduce the apparent long memory of realised volatility. The finding above says that the
correct exponent is not merely below $\tfrac12$ but very close to $0$: volatility is not
persistent, it is \emph{rough}. The resulting model, which we shall meet again in Chapter
\ref{ch:applications}, replaces the Heston dynamics of Subsection \ref{subsec:stochvol} by
\begin{equation}\label{eq:rfsv}
\sigma_t = \exp(X_t), \qquad \dd X_t = \nu\,\dd \BH_t - a\bigl(X_t - m\bigr)\dd t,
\end{equation}
with $a$ so small that, over any horizon of practical interest, $X$ is indistinguishable from
a fractional Brownian motion.\footnote{The mean-reversion speed is written $\alpha$ in
\cite{gatheraljaissonrosenbaum2018}; we use $a$, since $\alpha$ is reserved throughout this chapter
for the H\"older exponent fixed in \eqref{eq:alpharange}.}

For our purposes the decisive consequence is the one that concerns the object Chapter \ref{ch:finance}
was built to construct: the implied volatility surface. A model of the form
\eqref{eq:rfsv} generates an at-the-money skew whose term structure behaves like
\begin{equation}\label{eq:skewpowerlaw}
\psi(\tau) \;:=\; \biggl|\frac{\partial \ivol}{\partial \lmon}(\tau,\lmon)\biggr|_{\lmon=0} \;\sim\; \tau^{\Hurst - 1/2}
\qquad \text{as } \tau \downarrow 0 ,
\end{equation}
a power-law explosion at short maturities. Every model driven by a semimartingale volatility
factor --- Heston included, and by Theorem \ref{thm:fbmnotsemimartingale} that is precisely the
case $\Hurst = \tfrac12$ --- produces a skew that flattens as $\tau \downarrow 0$, and no choice of
parameters repairs this. The observed surface does not flatten. Roughness is what the market
implies, and $\eqref{eq:skewpowerlaw}$ is the cleanest statement of why.

\begin{remark}
The reader should note what \eqref{eq:rfsv} does and does not require. The equation is driven by
$\BH$ but is \emph{linear} in $X$, so it can be solved explicitly, by a pathwise
Riemann--Stieltjes argument, without any integration theory for $\BH$ at all. The difficulty
arrives one level up, when one wishes to price --- that is, to integrate $\sigma_t$, a nonlinear
functional of $\BH$, against the Brownian motion driving the asset price. That is a genuine
stochastic integral against a rough signal, and it is the problem this chapter exists to address.
\end{remark}

\subsection{What Fails, and Why}\label{subsec:whatfails}

Before constructing anything we should be clear that the obstruction is real, and in particular
that it is not removed by any sharpening of the classical estimates. There are three layers to
this, of increasing depth.

\medskip

\noindent\textbf{Young integration is out of reach.} Recall from Remark \ref{rem:criticalexponent}
that $\int g\,\dd f$ exists pathwise as a limit of Riemann sums whenever $f$ has finite
$p$-variation, $g$ finite $q$-variation, and $1/p + 1/q > 1$. In Hölder terms, if
$f \in \Holder{\alpha}$ and $g \in \Holder{\beta}$ the condition reads $\alpha + \beta > 1$. Now the
integrands we care about are of the form $g = F(f)$, which inherits the regularity $\alpha$ of $f$
itself, so the condition becomes $2\alpha > 1$, that is $\alpha > \tfrac12$. By Proposition
\ref{prop:fbmproperties} c) this holds for $\BH$ precisely when $\Hurst > \tfrac12$. Young's theory
therefore covers exactly the persistent case that the data reject, and none of the rough one.

\medskip

\noindent\textbf{Itô integration is out of reach.} By Theorem \ref{thm:fbmnotsemimartingale},
immediately and for all $\Hurst \neq \tfrac12$.

\medskip

\noindent\textbf{And the obstruction is not a matter of technique.} One might still hope that the
map $(f,g)\mapsto \int g\,\dd f$, manifestly well defined for smooth paths, extends by continuity
to the irregular ones --- that some cleverer topology on path space would do what the uniform
topology fails to do. The following example, due to Lyons, shows that the uniform topology cannot.

\begin{example}[The oscillation counterexample]\label{ex:lyonscounterexample}
For $n \in \N$ define the smooth planar paths
\begin{equation}\label{eq:lyonsexample}
x^{(n)}_t := \frac{1}{n}\sin\bigl(n^2 t\bigr), \qquad
y^{(n)}_t := \frac{1}{n}\cos\bigl(n^2 t\bigr), \qquad t \in [0,\Mat].
\end{equation}
Both converge to the zero path uniformly, since $|x^{(n)}|,|y^{(n)}| \leq 1/n$. Yet
\[
\int_0^t y^{(n)}_s\,\dd x^{(n)}_s = \int_0^t \frac{\cos(n^2s)}{n}\cdot n\cos(n^2 s)\,\dd s
= \int_0^t \cos^2(n^2 s)\,\dd s = \frac{t}{2} + \frac{\sin(2n^2t)}{4n^2} \;\longrightarrow\; \frac{t}{2},
\]
which is not the integral of the zero path against itself. The integral is therefore \emph{not}
continuous with respect to uniform convergence.
\end{example}

The picture behind the computation is worth holding on to, because it is the picture behind the
whole theory. The paths $\bigl(x^{(n)},y^{(n)}\bigr)$ trace circles of radius $1/n$ traversed $n^2$
times per unit time. The circles shrink to a point, which is why the paths converge uniformly to
zero; but the \emph{area} swept out is $\pi/n^2$ per revolution and there are $n^2/2\pi$
revolutions per unit time, so the total area accumulated per unit time is a constant, independent
of $n$. Uniform convergence sees the shrinking radius. It does not see the accumulating area. And
the integral, as the computation shows, depends on the area.

One might now hope for some other norm. There is none.

\begin{proposition}[Lyons]\label{prop:nonorm}
Let $\mathcal{W}$ denote the space of continuous paths $[0,\Mat]\to\R^d$ and let
$E^\infty \subseteq \mathcal{W}$ be the smooth paths. There is no norm $\lVert\cdot\rVert$ on
$E^\infty$ such that, simultaneously, the closure $E$ of $E^\infty$ carries full Wiener measure and
the map $(f,g)\mapsto\int g\,\dd f$ extends continuously from $E^\infty\times E^\infty$ to
$E\times E$. \qedopentext
\end{proposition}

A proof is given in \cite[Ch.\ 1]{frizhairer2020}. The statement is a genuine impossibility
theorem, not an admission of ignorance, and it forces a change of question. If no topology on
\emph{paths} makes the integral continuous, then the integral is not a function of the path alone,
and the object we should be topologising is not the path.

Example \ref{ex:lyonscounterexample} tells us exactly what is missing. What distinguishes
$\bigl(x^{(n)},y^{(n)}\bigr)$ from the zero path is the area it encloses --- equivalently, the
iterated integral $\int\!\!\int \dd x\,\dd y$. So we shall carry that quantity along as part of the
data, postulate it where it cannot be constructed, and topologise the enlarged object. That the
resulting theory is not merely a device but the \emph{correct} formulation is the content of
Theorem \ref{thm:ult}, and everything between here and there is the construction.

\medskip

One word on how the chapter is to be read, since its architecture is deliberate and easy to miss.
What follows is not a survey of an unrelated theory placed alongside Part \ref{part:part1}. Three
classical results --- Itô's formula and the Itô--Stratonovich correction, both obtained in Chapter
\ref{ch:foundations}, together with the Wong--Zakai approximation theorem, which belongs to the same
classical circle of ideas but was not needed there --- will be \emph{re-derived} here on
different foundations, and in each case we shall state explicitly what the new derivation costs in
hypotheses and what it buys in conclusions. That accounting is the real argument of this chapter.
The results themselves are not new, and the rough-path versions are in two of the three cases
strictly weaker than the classical ones; what is new is that they are obtained pathwise,
deterministically, and with a continuity statement attached. The coda closing the chapter collects
the three and draws the moral.

\section{Rough Paths}\label{sec:roughpaths}

\subsection{Tensor Spaces}\label{subsec:tensorspaces}

The iterated integrals we must carry are indexed by pairs of coordinates, and the algebraic home
for such objects is the tensor product. We recall only what is needed; the systematic treatment,
which requires free Lie algebras and the theory of the free nilpotent group, is developed in
\cite[Ch.\ 3]{geng2021} and \cite[Ch.\ 7]{frizvictoir2010}, and we shall not need it.

Throughout this chapter $V := \R^d$ with its standard basis $\{e_1,\dots,e_d\}$.

\begin{definition}[Tensor product]\label{def:tensorproduct}
The \textbf{tensor product} $V \otimes V$ is the vector space generated by the symbols
$v \otimes w$, for $v,w \in V$, subject to bilinearity:
\begin{equation}\label{eq:bilinearity}
(av_1 + v_2)\otimes w = a(v_1\otimes w) + v_2 \otimes w, \qquad
v \otimes (aw_1 + w_2) = a(v\otimes w_1) + v\otimes w_2 ,
\end{equation}
for all $a \in \R$ and $v,v_i,w,w_i \in V$.
\end{definition}

Concretely $V\otimes V$ has dimension $d^2$ with basis $\{e_i\otimes e_j\}_{1\leq i,j\leq d}$, so
that it may be identified with the space of $d\times d$ matrices, an element being written
$\xi = \sum_{i,j}\xi^{ij}\,e_i\otimes e_j$. The point of the notation is that the product is
\emph{not} commutative: $e_1\otimes e_2 \neq e_2\otimes e_1$. This is essential, since the two
iterated integrals $\int\!\!\int \dd x^1\dd x^2$ and $\int\!\!\int \dd x^2 \dd x^1$ are genuinely
different numbers, and it is exactly their difference that Example \ref{ex:lyonscounterexample}
detected.

\begin{definition}[Truncated tensor algebra]\label{def:tensoralgebra}
The \textbf{truncated tensor algebra of order two} over $V$ is
\begin{equation}\label{eq:tensoralgebra}
\Tens{2}(V) \;:=\; \R \oplus V \oplus (V\otimes V),
\end{equation}
whose elements we write $\xi = (\xi_0,\xi_1,\xi_2)$, equipped with the multiplication
\begin{equation}\label{eq:tensormultiplication}
(\xi \otimes \eta)_n \;:=\; \sum_{k=0}^{n} \xi_k \otimes \eta_{n-k}, \qquad n = 0,1,2,
\end{equation}
all terms of order higher than two being discarded. We write
$\Tone{2}(V) := \{\xi \in \Tens{2}(V) : \xi_0 = 1\}$. More generally, the truncation at order $N$ is
$\Tens{N}(V) := \bigoplus_{n=0}^{N} V^{\otimes n}$, with the same multiplication truncated at order
$N$, and the full tensor algebra --- in which no truncation is performed and elements are therefore
infinite sequences --- is $T((V)) := \prod_{n \geq 0} V^{\otimes n}$. Only $\Tens{2}$ is used until
Section \ref{sec:signature}.
\end{definition}

Written out, \eqref{eq:tensormultiplication} says
$\xi\otimes\eta = (\xi_0\eta_0,\ \xi_0\eta_1 + \xi_1\eta_0,\ \xi_0\eta_2 + \xi_1\otimes\eta_1 + \xi_2\eta_0)$.
The algebra is associative with unit $\mathbf{1} = (1,0,0)$, and although a general element has no
inverse, every element of $\Tone{2}(V)$ does, namely
\begin{equation}\label{eq:tensorinverse}
(1,\xi_1,\xi_2)^{-1} = \bigl(1,\ -\xi_1,\ \xi_1\otimes\xi_1 - \xi_2\bigr),
\end{equation}
as one checks by multiplying out. Thus $\Tone{2}(V)$ is a group, and this is the structure in which
the whole chapter takes place.

The motivation is the following computation, which the reader should regard as the definition of
everything that follows.

\begin{example}[The canonical lift of a smooth path]\label{ex:smoothlift}
Let $x : [0,\Mat]\to V$ be continuously differentiable. For $s \leq t$ set
\begin{equation}\label{eq:smoothlift}
X_{s,t} := x_t - x_s \in V, \qquad
\lev_{s,t} := \int_s^t \bigl(x_r - x_s\bigr)\otimes \dd x_r = \int_s^t\!\!\int_s^r \dd x_u \otimes \dd x_r \in V\otimes V,
\end{equation}
both integrals being ordinary Riemann--Stieltjes integrals. In coordinates,
$\lev^{ij}_{s,t} = \int_s^t (x^i_r - x^i_s)\,\dot x^j_r\,\dd r$. The pair
$\rpX_{s,t} := \bigl(1, X_{s,t}, \lev_{s,t}\bigr)$ takes values in $\Tone{2}(V)$.
\end{example}

Two features of \eqref{eq:smoothlift} will survive the passage to irregular paths, and one will not.
The two that survive are an algebraic identity, developed in the next subsection, and an analytic
bound. The one that does not is the following, which is worth isolating now.

\begin{proposition}[The symmetric part is determined]\label{prop:symmetricpart}
With the notation of Example \ref{ex:smoothlift},
\begin{equation}\label{eq:symmetricpart}
\Sym \lev_{s,t} \;:=\; \tfrac12\bigl(\lev_{s,t} + \lev_{s,t}\transp\bigr) \;=\; \tfrac12\,X_{s,t}\otimes X_{s,t}.
\end{equation}
\end{proposition}

\textbf{\textit{Proof:}} In coordinates the claim is
$\lev^{ij}_{s,t} + \lev^{ji}_{s,t} = X^i_{s,t}X^j_{s,t}$. Since $x$ is differentiable, the ordinary
product rule gives
\[
\frac{\dd}{\dd r}\Bigl[ \bigl(x^i_r - x^i_s\bigr)\bigl(x^j_r-x^j_s\bigr) \Bigr]
= \bigl(x^i_r-x^i_s\bigr)\dot x^j_r + \bigl(x^j_r-x^j_s\bigr)\dot x^i_r ,
\]
and integrating from $s$ to $t$, where the left-hand side vanishes at $r=s$, yields exactly the
claim. \qedfilledtext

\medskip

So for a smooth path the symmetric part of the iterated integral carries no information beyond the
increment itself: only the \emph{antisymmetric} part, the Lévy area
\begin{equation}\label{eq:levyarea}
A_{s,t} := \tfrac12\bigl(\lev_{s,t} - \lev_{s,t}\transp\bigr),
\end{equation}
is new. Geometrically $A^{ij}_{s,t}$ is the signed area enclosed between the chord joining
$(x^i_s,x^j_s)$ to $(x^i_t,x^j_t)$ and the arc of the path between them --- precisely the quantity
that Example \ref{ex:lyonscounterexample} showed to be invisible to uniform convergence. Whether
\eqref{eq:symmetricpart} continues to hold once $x$ is no longer differentiable is exactly the
distinction between the two classes of rough path introduced in Subsection \ref{subsec:typesofrp},
and it is the precise mathematical location of the difference between Itô and Stratonovich calculus.

\subsection{The Multiplicative Property}\label{subsec:multiplicative}

The algebraic identity announced above is the following, and it is the axiom on which the entire
theory rests.

\begin{proposition}[Chen's relation for smooth paths]\label{prop:chensmooth}
With the notation of Example \ref{ex:smoothlift}, for all $s \leq u \leq t$,
\begin{equation}\label{eq:chen}
\lev_{s,t} - \lev_{s,u} - \lev_{u,t} \;=\; X_{s,u}\otimes X_{u,t},
\end{equation}
equivalently $\rpX_{s,t} = \rpX_{s,u}\otimes\rpX_{u,t}$ in the group $\Tone{2}(V)$.
\end{proposition}

\textbf{\textit{Proof:}} Split the outer integral in \eqref{eq:smoothlift} at $u$:
\[
\lev_{s,t} = \int_s^u (x_r - x_s)\otimes \dd x_r + \int_u^t (x_r-x_s)\otimes \dd x_r
= \lev_{s,u} + \int_u^t \bigl[(x_r - x_u) + (x_u - x_s)\bigr]\otimes \dd x_r .
\]
The first term of the bracket integrates to $\lev_{u,t}$. The second is constant in $r$, so it
factors out of the integral, leaving $X_{s,u}\otimes\int_u^t \dd x_r = X_{s,u}\otimes X_{u,t}$. That
the resulting three identities --- one at each level --- are precisely the statement
$\rpX_{s,t} = \rpX_{s,u}\otimes\rpX_{u,t}$ is immediate from \eqref{eq:tensormultiplication}.
\qedfilledtext

\medskip

Relation \eqref{eq:chen} is worth reading twice. At the first level it is the triviality
$X_{s,t} = X_{s,u} + X_{u,t}$, the additivity of increments. At the second level it is a genuine
constraint, and it says that $\lev$ is \emph{not} additive: concatenating two intervals produces a
cross-term $X_{s,u}\otimes X_{u,t}$ which is the area contributed by the interaction of the two
pieces. The relation is thus the algebraic residue of the additivity of integration, and it is what
makes $\lev$ a coherent object rather than an arbitrary two-parameter function.

\begin{notationbox}
Throughout, $\lev$ is a function of \emph{two} time variables and is never a path. There is no
``$\lev_t$''. This is not a notational accident: by \eqref{eq:chen} the increments of $\lev$ do not
telescope, so no single-variable function generates it. The first level, by contrast, may be
regarded either way, and we write $X_{s,t} = x_t - x_s$ when convenient.
\end{notationbox}

We can now give the central definition. From this point onwards we fix
\begin{equation}\label{eq:alpharange}
\alpha \in \bigl(\tfrac13, \tfrac12\bigr],
\end{equation}
a restriction whose necessity will become visible in Theorem \ref{thm:gubinelli} and which we
discuss in Remark \ref{rem:whyonethird}.

A word on why the exponent is a Hölder one. Lyons' original theory \cite{lyons1998}, and the
systematic treatment of \cite{frizvictoir2010}, are formulated in $p$-variation --- the scale of
Definition \ref{def:pvariation}, in which Proposition \ref{prop:fbmproperties} d) and Remark
\ref{rem:criticalexponent} are also stated, so that a reader arriving here expects it. The two
settings are essentially interchangeable, the dictionary being $p = 1/\alpha$: a path of finite
$p$-variation is $\alpha$-Hölder after reparametrisation by its own variation, and every
$\alpha$-Hölder path has finite $1/\alpha$-variation. We work in Hölder norms because they make the
estimates of Subsection \ref{subsec:roughintegration} additive in the exponent and therefore easier
to read, and because \eqref{eq:alpharange} then corresponds to $p \in [2,3)$ --- exactly the range
in which two levels suffice.

\begin{definition}[Rough path]\label{def:roughpath}
An \textbf{$\alpha$-Hölder rough path} over $V$ is a pair $\rpX = (X,\lev)$ consisting of
$X \in \Holder{\alpha}\bigl([0,\Mat],V\bigr)$ and a two-parameter function
$\lev : \{0\leq s\leq t\leq \Mat\} \to V\otimes V$ such that:
\begin{itemize}
    \item[i)] (\textbf{Chen's relation}) for all $s\leq u\leq t$,
    \begin{equation}\label{eq:chenaxiom}
    \lev_{s,t} - \lev_{s,u} - \lev_{u,t} = X_{s,u}\otimes X_{u,t};
    \end{equation}
    \item[ii)] (\textbf{analytic bound}) the seminorms
    \begin{equation}\label{eq:rpnorms}
    \lVert X\rVert_{\alpha} := \sup_{s\neq t}\frac{|X_{s,t}|}{|t-s|^{\alpha}}, \qquad
    \lVert \lev\rVert_{2\alpha} := \sup_{s\neq t}\frac{|\lev_{s,t}|}{|t-s|^{2\alpha}}
    \end{equation}
    are both finite.
\end{itemize}
We write $\rpX \in \RP{\alpha}\bigl([0,\Mat],V\bigr)$, abbreviated $\RP{\alpha}$, and set
$\vertiii{\rpX}_{\alpha} := \lVert X\rVert_\alpha + \lVert\lev\rVert_{2\alpha}^{1/2}$.
\end{definition}

The exponent $2\alpha$ in \eqref{eq:rpnorms} is not a convention but a necessity: $\lev$ is
formally a double integral of increments of size $|t-s|^\alpha$, so any object deserving the name
must scale that way. Note also that the two conditions pull in opposite directions --- i) is
algebraic and rigid, ii) is analytic and merely quantitative --- and that neither alone would
suffice.

\begin{definition}[The rough path metric]\label{def:rpmetric}
For $\rpX = (X,\lev)$ and $\widetilde{\rpX} = (\widetilde X, \widetilde\lev)$ in $\RP{\alpha}$ set
\begin{equation}\label{eq:rpmetric}
\rpdist\bigl(\rpX,\widetilde{\rpX}\bigr) \;:=\; \bigl\lVert X - \widetilde X\bigr\rVert_{\alpha} \;+\; \bigl\lVert \lev - \widetilde\lev\bigr\rVert_{2\alpha}.
\end{equation}
\end{definition}

The second level enters \eqref{eq:rpmetric} to the first power, whereas it entered
$\vertiii{\rpX}_\alpha$ in Definition \ref{def:roughpath} under a square root, so that
$\rpdist(\rpX,\mathbf{0}) \neq \vertiii{\rpX}_\alpha$. This is not a misprint. The metric is
\emph{inhomogeneous} and is the one in which the estimates of Subsection \ref{subsec:ult} are
linear; the norm is \emph{homogeneous}, meaning that it scales correctly under the dilation
$(X,\lev)\mapsto(\lambda X,\lambda^2\lev)$ which the exponent $2\alpha$ of \eqref{eq:rpnorms}
forces. Both are standard, and each is used where its own scaling behaviour is wanted.

\begin{remark}[$\RP{\alpha}$ is not a vector space]\label{rem:notvectorspace}
The set of rough paths is closed under neither addition nor scalar multiplication, since
\eqref{eq:chenaxiom} is quadratic in $X$: if $\rpX$ and $\widetilde\rpX$ satisfy it, their sum does
not. It is a complete metric space under $\rpdist$, and a closed subset of the Banach space
$\Holder{\alpha}\oplus \Holder{2\alpha}_2$ --- the subscript $2$ denoting functions of \emph{two}
time variables, as $\lev$ is --- in which it naturally sits, but it is neither a linear
space nor separable. The correct picture is a nonlinear ``fibre bundle'': over each $\alpha$-Hölder
path $X$ sits the set of admissible second levels, which we now describe. See \cite[Ch.\ 2]{frizhairer2020}.
\end{remark}

\begin{proposition}[The lift is not unique]\label{prop:liftnotunique}
Let $\rpX = (X,\lev) \in \RP{\alpha}$ and let $f : [0,\Mat] \to V\otimes V$ belong to
$\Holder{2\alpha}$. Then $\widetilde\rpX := \bigl(X,\ \lev + \increment{f}\bigr)$, where
$\increment{f}_{s,t} := f_t - f_s$, is again an element of $\RP{\alpha}$. Conversely, if
$(X,\lev)$ and $(X,\widetilde\lev)$ both lie in $\RP{\alpha}$, then
$\widetilde\lev - \lev = \increment{f}$ for some $f \in \Holder{2\alpha}$.
\end{proposition}

\textbf{\textit{Proof:}} For the first claim, note that increments telescope: $\increment{f}_{s,t} -
\increment{f}_{s,u} - \increment{f}_{u,t} = 0$, so adding $\increment{f}$ to $\lev$ leaves
\eqref{eq:chenaxiom} unchanged, while $\lVert\increment{f}\rVert_{2\alpha} = \lVert f\rVert_{2\alpha}
< \infty$ preserves the bound. Conversely, set $\Delta_{s,t} := \widetilde\lev_{s,t} - \lev_{s,t}$.
Subtracting the two instances of \eqref{eq:chenaxiom}, which share the same right-hand side because
the first levels agree, gives $\Delta_{s,t} = \Delta_{s,u} + \Delta_{u,t}$ for all $s\leq u\leq t$.
Taking $s=u$ shows $\Delta_{s,s}=0$, and then $f_t := \Delta_{0,t}$ satisfies
$\increment{f}_{s,t} = \Delta_{0,t}-\Delta_{0,s} = \Delta_{s,t}$, with
$\lVert f\rVert_{2\alpha} = \lVert\Delta\rVert_{2\alpha} < \infty$ by the triangle inequality.
\qedfilledtext

\medskip

This is the single most important structural fact in the subject, and it should be dwelt on.
\textbf{A rough path is not a path.} The second level is genuine additional data, undetermined by
the trajectory, and the space of admissible choices over a fixed $X$ is an affine copy of
$\Holder{2\alpha}$. The theory does not \emph{discover} $\lev$; it takes $\lev$ as given and
proceeds. Different choices are different rough paths, and --- as Remark \ref{rem:integraldepends}
will make precise --- they yield different integrals and different solutions to differential
equations. Far from being a defect, this is what allows a single formalism to contain both the Itô
and the Stratonovich calculus of Chapter \ref{ch:foundations} as special cases, distinguished by nothing
other than the choice of $\lev$.

\subsection{Types of Rough Paths}\label{subsec:typesofrp}

Proposition \ref{prop:symmetricpart} held automatically for smooth paths. Promoting it to a
definition isolates the class on which ordinary calculus survives.

\begin{definition}[Geometric rough paths]\label{def:geometric}
A rough path $\rpX = (X,\lev) \in \RP{\alpha}$ is \textbf{weakly geometric} if
\begin{equation}\label{eq:weaklygeometric}
\Sym\lev_{s,t} = \tfrac12\, X_{s,t}\otimes X_{s,t} \qquad \text{for all } s \leq t ,
\end{equation}
and \textbf{geometric} if it is the limit, in the metric $\rpdist$, of canonical lifts of smooth
paths in the sense of Example \ref{ex:smoothlift}. We write $\RPwg{\alpha}$ and $\RPg{\alpha}$
respectively.
\end{definition}

\begin{proposition}\label{prop:geometricinclusion}
$\RPg{\alpha} \subseteq \RPwg{\alpha}$.
\end{proposition}

\textbf{\textit{Proof:}} Each canonical lift satisfies \eqref{eq:weaklygeometric} by Proposition
\ref{prop:symmetricpart}. The condition is a closed one: if $\rpdist(\rpX^{(n)},\rpX)\to0$ then in
particular $\lev^{(n)}_{s,t}\to\lev_{s,t}$ and $X^{(n)}_{s,t}\to X_{s,t}$ for each fixed pair $s,t$,
and both sides of \eqref{eq:weaklygeometric} pass to the limit. \qedfilledtext

\medskip

The inclusion is strict, but only just, and for our purposes harmlessly: a weakly geometric
$\alpha$-Hölder rough path is a geometric $\beta$-Hölder rough path for every $\beta<\alpha$, a
result of Friz and Victoir \cite[Ch.\ 8]{frizvictoir2010}. Since our hypotheses always tolerate a
slightly smaller exponent, we shall use the two notions interchangeably and say simply
\emph{geometric}, flagging the distinction only where it does work.

The content of \eqref{eq:weaklygeometric} is best seen through what it makes true. If $\rpX$ is
geometric then the chain rule takes its classical form, with no second-order correction: geometric
rough paths are exactly those for which the calculus is ordinary calculus. We are not yet in a
position to say this precisely --- the statement involves an integral we have not built --- but it
is proved as Proposition \ref{prop:roughchainrule} below, and the reader may wish to read that
result as the true definition of ``geometric''. Non-geometric rough paths are those for which the
calculus is not ordinary --- and the canonical example is one we have already met at length.

\begin{example}[Brownian motion as a rough path]\label{ex:brownianroughpath}
Let $W$ be a $d$-dimensional Brownian motion. Define, for $s\leq t$,
\begin{equation}\label{eq:itolift}
\levW^{\mathrm{It\hat o}}_{s,t} := \int_s^t \bigl(W_r - W_s\bigr)\otimes \dd W_r ,
\end{equation}
the integral being the Itô integral of Definition \ref{def:stochastic_integral}, and
\begin{equation}\label{eq:stratlift}
\levW^{\mathrm{Strat}}_{s,t} := \levW^{\mathrm{It\hat o}}_{s,t} + \tfrac12 (t-s)\,\mathrm{Id},
\end{equation}
where $\mathrm{Id} = \sum_i e_i \otimes e_i$. Then, for every $\alpha \in (\tfrac13,\tfrac12)$, both
$\mathbf{W}^{\mathrm{It\hat o}} := (W,\levW^{\mathrm{It\hat o}})$ and
$\mathbf{W}^{\mathrm{Strat}} := (W,\levW^{\mathrm{Strat}})$ are almost surely rough paths in the
sense of Definition \ref{def:roughpath}. The second is geometric; the first is not.
\end{example}

\textbf{\textit{Proof:}} Chen's relation \eqref{eq:chenaxiom} for $\levW^{\mathrm{It\hat o}}$ is the
computation of Proposition \ref{prop:chensmooth} verbatim, the splitting of the outer integral at
$u$ and the extraction of the constant $W_u - W_s$ being valid for the Itô integral by linearity.
For \eqref{eq:stratlift} it holds because the added term $\tfrac12(t-s)\mathrm{Id}$ is an increment
and therefore telescopes, exactly as in Proposition \ref{prop:liftnotunique}. The analytic bound
\eqref{eq:rpnorms} follows from the Itô isometry (Theorem \ref{thm:itoisometry}), which gives
$\E\bigl[|\levW^{\mathrm{It\hat o}}_{s,t}|^2\bigr] \leq C|t-s|^2$, hence
$\E\bigl[|\levW_{s,t}|^{2m}\bigr] \leq C_m |t-s|^{2m}$ by Gaussian hypercontractivity, followed by a
two-parameter version of Kolmogorov's criterion; the details are in \cite[Ch.\ 3]{frizhairer2020}.

For the geometric assertion, apply the multidimensional Itô rule (Theorem \ref{thm:itomulti}) to the
product $W^i_{s,t}W^j_{s,t}$, using the Itô multiplication table of Chapter \ref{ch:foundations}:
\[
W^i_{s,t}W^j_{s,t} = \int_s^t W^i_{s,r}\,\dd W^j_r + \int_s^t W^j_{s,r}\,\dd W^i_r + \delta_{ij}(t-s),
\]
the last term being $\langle W^i, W^j\rangle_{s,t}$ by Proposition \ref{prop:qvcoincide}. In tensor
notation this reads $2\,\Sym\levW^{\mathrm{It\hat o}}_{s,t} = W_{s,t}\otimes W_{s,t} - (t-s)\mathrm{Id}$,
so \eqref{eq:weaklygeometric} fails for the Itô lift by exactly $-\tfrac12(t-s)\mathrm{Id}$ --- and
holds for the Stratonovich lift, which was defined by adding precisely that discrepancy back.
\qedfilledtext

\medskip

\begin{remark}[The law of the Lévy area]\label{rem:levyarealaw}
Since $\levW$ has just been constructed, it is worth recording what its distribution is. For $d=2$
the antisymmetric part \eqref{eq:levyarea} is the scalar
$A_{s,t} = \tfrac12\bigl(\lev^{12}_{s,t}-\lev^{21}_{s,t}\bigr)$, the signed area swept by the planar
Brownian path against its chord, and Lévy's formula gives its conditional characteristic function
\begin{equation}\label{eq:levyareaformula}
\E\Bigl[e^{iuA_{s,t}} \,\Big|\, W_{s,t} = z\Bigr]
\;=\; \frac{u(t-s)/2}{\sinh\bigl(u(t-s)/2\bigr)}\,
\exp\left( \frac{|z|^2}{2(t-s)}\Bigl[1 - \tfrac{u(t-s)}{2}\coth\tfrac{u(t-s)}{2}\Bigr]\right).
\end{equation}
Unconditionally, $A_{0,t}$ is distributed as $\tfrac{t}{2}\,C$ with $C$ having characteristic
function $1/\cosh(u/2)$. Note that the law is \emph{not} Gaussian, which is the first sign that the
second level is a genuinely new object rather than a function of the first; the derivation, via the
Cameron--Martin formula for a quadratic functional of Brownian motion, is in
\cite[Ch.\ 13]{frizvictoir2010}. We do not use \eqref{eq:levyareaformula}, but it is what one would
simulate from if constructing $\levW$ directly rather than by discretisation.
\end{remark}

The reader should pause over what has just happened. The Itô--Stratonovich correction of Chapter
\ref{ch:foundations}, which appeared there as a consequence of the left-endpoint convention in the
construction of the integral, has reappeared here as the difference between two \emph{elements of
the same space} $\RP{\alpha}$, lying over the same first level and differing by the increment of
$f_t = \tfrac12 t\,\mathrm{Id} \in \Holder{1} \subset \Holder{2\alpha}$. Proposition
\ref{prop:liftnotunique} predicted that such pairs exist; Example \ref{ex:brownianroughpath} says
that the two calculi of Part \ref{part:part1} are an instance. The choice between them is not
forced by the trajectory. It never was.

We turn to the process that motivated the chapter.

\begin{theorem}[Coutin--Qian]\label{thm:fbmlift}
Let $\BH$ be a $d$-dimensional fractional Brownian motion with $\Hurst \in (\tfrac13,\tfrac12]$ and
independent components, and let $W^{\Hurst,(n)}$ denote its piecewise-linear interpolation over the
dyadic partition of mesh $2^{-n}$, canonically lifted as in Example \ref{ex:smoothlift}. Then, for
every $\alpha < \Hurst$, the lifts converge in probability in the metric $\rpdist$ to a limit
$\mathbf{B}^{\Hurst} = \bigl(\BH,\lev^{\Hurst}\bigr) \in \RPg{\alpha}$, which is therefore a geometric
rough path. \qedopentext
\end{theorem}

The boldface letter is $\mathbf{B}$ and not $\mathbf{W}$ deliberately: the first level of
$\mathbf{B}^{\Hurst}$ is the fractional Brownian motion $\BH$, but $\mathbf{W}^{\mathrm{It\hat o}}$
and $\mathbf{W}^{\mathrm{Strat}}$ of Example \ref{ex:brownianroughpath} are lifts of an ordinary
Brownian motion, and the two families are never the same object.

The result is due to Coutin and Qian \cite{coutinqian2002}; a modern proof, via the
$\varrho$-variation of the covariance \eqref{eq:fbmcovariance}, is given in
\cite[Ch.\ 10]{frizhairer2020}. The mechanism is instructive: one shows that
$R_\Hurst$ has finite two-dimensional $\varrho$-variation with $\varrho = 1/(2\Hurst)$, and the
construction requires $\varrho < \tfrac32$, which is exactly $\Hurst > \tfrac13$.

\begin{remark}[The empirical exponent lies outside]\label{rem:H01outside}
Theorem \ref{thm:fbmlift} covers $\Hurst > \tfrac13$; by admitting a third iterated integral --- a
level we have deliberately excluded from Definition \ref{def:roughpath} --- the range extends to
$\Hurst > \tfrac14$, and no further, since for $\Hurst \leq \tfrac14$ the piecewise-linear lifts
fail to converge at all. The exponent $\Hurst \approx 0.1$ reported in Subsection \ref{subsec:roughvol} therefore lies
outside the reach of the canonical construction, at any number of levels. The precise statement
worth remembering is not that no lift exists --- non-canonical liftings of $\BH$ below $\tfrac14$
have been constructed by other means, the Lyons--Victoir extension theorem
\cite{lyonsvictoir2007} guaranteeing by a Hahn--Banach argument that \emph{every} $\alpha$-Hölder
path admits at least one lift, at the cost of the axiom of choice and with no canonical selection
among them --- but that no lift exists which is the limit of the
piecewise-linear approximations. That is the property one actually needs: a lift that no
discretisation converges to is a lift that cannot be simulated, cannot be estimated from data, and
carries no guarantee that a numerical scheme is approximating it rather than something else. The
obstruction is thus not an absence of objects but an absence of \emph{canonical} ones.

This is not a gap in the exposition but a real boundary, and it would be dishonest to leave it
unmarked. Section \ref{sec:beyond} says what lies beyond it. Two remarks temper the difficulty in
the meantime. First, the model \eqref{eq:rfsv} is linear in $X$ and needs no integration theory for
$\BH$; the rough integral is required for the \emph{pricing} equation, where the integrator is the
Brownian motion driving the asset and $\Hurst$ does not enter the regularity of the integrator at
all. Second, and more importantly for Chapter \ref{ch:applications}, the signature of Section
\ref{sec:signature} is a transform of an \emph{observed} data stream, and observed data streams are
piecewise linear. No lifting theorem is needed to compute it, and no value of $\Hurst$ obstructs it.
\end{remark}

\section{Fundamental Results}\label{sec:fundamentalresults}

\subsection{Integration}\label{subsec:roughintegration}

We now construct the integral. The strategy has three steps: an abstract lemma which converts local
estimates into a global object; a class of admissible integrands; and the theorem that puts them
together.

\medskip

\noindent\textbf{Step one: sewing.} Every construction of an integral in this dissertation has had
the same shape --- propose a Riemann sum, show it converges. The following lemma, due in this form
to Gubinelli, extracts the mechanism once and for all.

\begin{lemma}[Sewing Lemma]\label{lem:sewing}
Let $\beta > 1$ and let $\Xi : \{0 \leq s \leq t \leq \Mat\} \to \R^m$ be continuous with
$\Xi_{t,t}=0$ and
\begin{equation}\label{eq:sewinghypothesis}
\bigl| \increment{\Xi}_{s,u,t} \bigr| \;\leq\; C\,|t-s|^{\beta}, \qquad \text{where }
\increment{\Xi}_{s,u,t} := \Xi_{s,t} - \Xi_{s,u} - \Xi_{u,t},
\end{equation}
for all $s\leq u\leq t$. Then there exists a unique function $\mathcal{I}\Xi$, of the same two
variables, which is \emph{additive} --- that is, $(\mathcal{I}\Xi)_{s,t} = (\mathcal{I}\Xi)_{s,u} +
(\mathcal{I}\Xi)_{u,t}$ for all $s\leq u\leq t$ --- and satisfies
\begin{equation}\label{eq:sewingconclusion}
\bigl| (\mathcal{I}\Xi)_{s,t} - \Xi_{s,t} \bigr| \;\leq\; \frac{C}{1 - 2^{\,1-\beta}}\,|t-s|^{\beta}.
\end{equation}
Moreover $(\mathcal{I}\Xi)_{s,t} = \lim_{|\mathcal{P}|\to0}\sum_{[u,v]\in\mathcal{P}} \Xi_{u,v}$,
the limit being over partitions $\mathcal{P}$ of $[s,t]$ with mesh tending to zero.
\end{lemma}

\textbf{\textit{Proof (uniqueness):}} Suppose $A$ and $\widetilde A$ are both additive and both
satisfy \eqref{eq:sewingconclusion}, with constants $C_1$ and $C_2$. Then $D := A - \widetilde A$ is
additive and $|D_{s,t}| \leq (C_1+C_2)|t-s|^\beta =: K|t-s|^\beta$. Fix $s<t$ and let
$\mathcal{P}_n$ be the uniform partition of $[s,t]$ into $n$ pieces. Additivity gives
\[
|D_{s,t}| = \Bigl|\sum_{[u,v]\in\mathcal{P}_n} D_{u,v}\Bigr| \leq \sum_{[u,v]\in\mathcal{P}_n} K\,|v-u|^\beta
= n \cdot K \Bigl(\frac{t-s}{n}\Bigr)^{\beta} = K\,|t-s|^{\beta}\,n^{1-\beta}.
\]
Since $\beta>1$ the right-hand side tends to $0$ as $n\to\infty$, so $D \equiv 0$. The existence
half, which proceeds by dyadic refinement, is proved in Appendix \ref{app:proof:sewing}.
\qedfilledtext

\medskip

The hypothesis \eqref{eq:sewinghypothesis} is the exact quantitative statement of ``$\Xi$ is almost
additive'', and the conclusion is that an almost additive function is uniformly close to a unique
genuinely additive one. The strictness of $\beta>1$ is what the counting argument above consumes,
and it is not improvable: at $\beta=1$ the factor $n^{1-\beta}$ is $1$ and the argument collapses.
Everything in this section is an application of the lemma with a specific $\Xi$, and the whole
difficulty is arranging for $\beta>1$ to hold.

\medskip

\noindent\textbf{Step two: admissible integrands.} We cannot integrate an arbitrary path against
$\rpX$; the integrand must be compatible with the integrator in a sense that Chapter \ref{ch:foundations}
would have called adaptedness, but which is here entirely deterministic.

\begin{definition}[Controlled rough path]\label{def:controlled}
Let $\rpX = (X,\lev)\in\RP{\alpha}$. A pair $(Y,\Gub)$ of paths
$Y \in \Holder{\alpha}([0,\Mat],W)$ and
$\Gub \in \Holder{\alpha}\bigl([0,\Mat],\Lloc{V,W}\bigr)$ is a path \textbf{controlled by} $\rpX$ if
the remainder defined by
\begin{equation}\label{eq:controlleddef}
R^{Y}_{s,t} \;:=\; Y_{s,t} - \Gub_s\,X_{s,t}, \qquad Y_{s,t} := Y_t - Y_s,
\end{equation}
satisfies $\lVert R^Y\rVert_{2\alpha} < \infty$. We write $(Y,\Gub) \in \Ctrl{X}$ and call $\Gub$
the \textbf{Gubinelli derivative} of $Y$ with respect to $X$.
\end{definition}

Equation \eqref{eq:controlleddef} is a first-order Taylor expansion of $Y$ along $X$: the path $Y$
is required to look, locally, like a linear image of $X$, with an error of order $2\alpha$ rather
than the $\alpha$ that mere Hölder continuity would give. The pair $\bigl(\Ctrl{X}, \lVert\cdot\rVert\bigr)$
with $\lVert (Y,\Gub)\rVert_{X,2\alpha} := \lVert \Gub\rVert_\alpha + \lVert R^Y\rVert_{2\alpha}$ is
a Banach space once the initial values are fixed \cite[Ch.\ 4]{frizhairer2020}.

\begin{remark}[$\Gub$ is data, not a derivative]\label{rem:gubinellinotunique}
The notation $\Gub$ is suggestive but must not be taken literally: in general $\Gub$ is \emph{not}
determined by $Y$ and $X$. If $X$ happens to be smooth, say $X_t = t$, then any $\Gub$ whatsoever
satisfies \eqref{eq:controlleddef} with a remainder of the required order, since
$\Gub_s X_{s,t} = O(|t-s|)$ is itself $2\alpha$-small when $2\alpha \leq 1$. Uniqueness holds only
when $X$ is genuinely irregular --- ``truly rough'' in the terminology of \cite[Ch.\ 6]{frizhairer2020}
--- which is the case almost surely for Brownian motion and for $\BH$. Like the second level $\lev$
itself, $\Gub$ is part of the given data. The theory is consistently one in which one carries more
information than the classical objects contain.
\end{remark}

\begin{example}\label{ex:controlledexamples}
i) The integrator controls itself: $(X, \mathrm{Id}) \in \Ctrl{X}$, with $R^X \equiv 0$.

ii) If $F \in \mathcal{C}^2_b(W,\widetilde W)$ and $(Y,\Gub) \in \Ctrl{X}$, then
$\bigl(F(Y),\,DF(Y)\Gub\bigr) \in \Ctrl{X}$. Indeed, Taylor's theorem with integral remainder gives
\[
F(Y_t)-F(Y_s) = DF(Y_s)Y_{s,t} + O\bigl(|Y_{s,t}|^2\bigr)
= DF(Y_s)\Gub_s X_{s,t} + DF(Y_s)R^Y_{s,t} + O\bigl(|t-s|^{2\alpha}\bigr),
\]
and the last two terms are both of order $2\alpha$. This closes the class under the nonlinearities
that appear in differential equations, and is the reason Definition \ref{def:controlled} is the
right one.
\end{example}

\medskip

\noindent\textbf{Step three: the integral.} The naïve Riemann sum $\sum Y_u X_{u,v}$ does not
converge: its defect is of order $|t-s|^{2\alpha}$, and $2\alpha \leq 1$, so Lemma \ref{lem:sewing}
does not apply. The remedy is to add a second-order term, and the expansion \eqref{eq:controlleddef}
tells us precisely which one.

\begin{theorem}[Gubinelli]\label{thm:gubinelli}
Let $\alpha \in (\tfrac13,\tfrac12]$, let $\rpX = (X,\lev)\in\RP{\alpha}$ and let
$(Y,\Gub) \in \Ctrl{X}$. Then the \textbf{compensated Riemann sums}
\begin{equation}\label{eq:roughintegral}
\int_s^t Y_r\,\dd\rpX_r \;:=\; \lim_{|\mathcal{P}|\to0} \sum_{[u,v]\in\mathcal{P}} \Bigl( Y_u\,X_{u,v} \;+\; \Gub_u\,\lev_{u,v} \Bigr)
\end{equation}
converge, and the limit satisfies
\begin{equation}\label{eq:roughintegralestimate}
\biggl| \int_s^t Y_r\,\dd\rpX_r - Y_s X_{s,t} - \Gub_s \lev_{s,t} \biggr|
\;\leq\; C\Bigl( \lVert X\rVert_\alpha \lVert R^Y\rVert_{2\alpha} + \lVert \lev\rVert_{2\alpha}\lVert \Gub\rVert_\alpha \Bigr)\,|t-s|^{3\alpha},
\end{equation}
for a constant $C = C(\alpha)$. Moreover the indefinite integral is itself controlled by $\rpX$:
\begin{equation}\label{eq:integraliscontrolled}
\Bigl( \int_0^{\cdot} Y_r\,\dd\rpX_r,\ Y \Bigr) \;\in\; \Ctrl{X}.
\end{equation}
\end{theorem}

The proof, an application of Lemma \ref{lem:sewing} to
$\Xi_{s,t} := Y_sX_{s,t} + \Gub_s\lev_{s,t}$, is given in Appendix \ref{app:proof:gubinelli}.

\begin{remark}[Where $\alpha>\tfrac13$ comes from]\label{rem:whyonethird}
The exponent $3\alpha$ in \eqref{eq:roughintegralestimate} is the whole story. Computing
$\increment{\Xi}_{s,u,t}$ for the choice above, and using Chen's relation \eqref{eq:chenaxiom} to
cancel the leading terms --- this is the computation in the appendix, and it is the only place
where \eqref{eq:chenaxiom} is used --- one finds
\begin{equation}\label{eq:xidefect}
-\increment{\Xi}_{s,u,t} = R^{Y}_{s,u}X_{u,t} + \bigl(\Gub_u - \Gub_s\bigr)\lev_{u,t},
\end{equation}
whose two terms are of order $2\alpha + \alpha$ and $\alpha + 2\alpha$ respectively. So
$\beta = 3\alpha$, and Lemma \ref{lem:sewing} demands $\beta > 1$, that is $\alpha > \tfrac13$.
Below that threshold the second level no longer suffices and a third iterated integral must be
carried; below $\tfrac14$, a fourth; and in general $\lfloor 1/\alpha\rfloor$ levels are required.
We have restricted to \eqref{eq:alpharange} precisely to avoid the algebra that the general case
demands, at no cost to the ideas.
\end{remark}

\begin{theorem}[Stability of the integral]\label{thm:integralstability}
Fix $M<\infty$ and consider $\rpX,\widetilde\rpX \in \RP{\alpha}$ and $(Y,\Gub)\in\Ctrl{X}$,
$(\widetilde Y,\widetilde\Gub)\in\Ctrl{\widetilde X}$, all of norm at most $M$. Then, with
$C = C(M,\alpha)$,
\begin{equation}\label{eq:integralstability}
\Bigl\lVert \int_0^{\cdot} Y\,\dd\rpX - \int_0^{\cdot} \widetilde Y\,\dd\widetilde\rpX \Bigr\rVert_{\alpha}
\;\leq\; C\Bigl( \rpdist\bigl(\rpX,\widetilde\rpX\bigr) + \bigl|Y_0 - \widetilde Y_0\bigr| + \bigl|\Gub_0 - \widetilde\Gub_0\bigr| + \lVert Y,\Gub;\widetilde Y,\widetilde \Gub\rVert_{X,\widetilde X,2\alpha} \Bigr),
\end{equation}
where the last term is the distance between two controlled paths lying over \emph{different}
rough paths,
\begin{equation}\label{eq:controlleddistance}
\lVert Y,\Gub;\widetilde Y,\widetilde\Gub\rVert_{X,\widetilde X,2\alpha}
\;:=\; \bigl\lVert \Gub - \widetilde\Gub \bigr\rVert_{\alpha} \;+\; \bigl\lVert R^{Y} - R^{\widetilde Y}\bigr\rVert_{2\alpha},
\end{equation}
with $R^Y_{s,t} = Y_{s,t} - \Gub_sX_{s,t}$ and
$R^{\widetilde Y}_{s,t} = \widetilde Y_{s,t} - \widetilde\Gub_s\widetilde X_{s,t}$ as in
\eqref{eq:controlleddef}. \qedopentext
\end{theorem}

The quantity \eqref{eq:controlleddistance} requires a word, since $(Y,\Gub)$ and
$(\widetilde Y,\widetilde\Gub)$ live in the \emph{different} Banach spaces $\Ctrl{X}$ and
$\Ctrl{\widetilde X}$ and cannot simply be subtracted. What is compared instead is the pair of
Gubinelli derivatives, which live in a common space, together with the pair of remainders, which
also do; the two rough paths enter only through the definitions of $R^Y$ and $R^{\widetilde Y}$.
This is the standard device, and it is the reason the estimate carries $\rpdist(\rpX,\widetilde\rpX)$
as a separate term rather than absorbing it.

This is proved in \cite[Thm.\ 4.17]{frizhairer2020} by re-running the proof of Theorem
\ref{thm:gubinelli} on differences. We record it because it, and not Theorem \ref{thm:gubinelli}
itself, is what makes the continuity statement of Subsection \ref{subsec:ult} possible.

The construction would be worthless if it disagreed with the theory of Part \ref{part:part1} where
both apply. It does not.

\begin{proposition}[Consistency]\label{prop:consistency}
Let $W$ be a Brownian motion and let $F \in \mathcal{C}^3_b$. Then, almost surely,
\begin{equation}\label{eq:consistency}
\int_0^t F(W_r)\,\dd\mathbf{W}^{\mathrm{It\hat o}}_r = \int_0^t F(W_r)\,\dd W_r,
\qquad
\int_0^t F(W_r)\,\dd\mathbf{W}^{\mathrm{Strat}}_r = \int_0^t F(W_r)\circ \dd W_r,
\end{equation}
the right-hand sides being the Itô integral of Definition \ref{def:stochastic_integral} and the
Stratonovich integral of Definition \ref{def:stratonovich} respectively. If $X$ is smooth, the rough integral against its canonical lift
is the Riemann--Stieltjes integral. \qedopentext
\end{proposition}

See \cite[Ch.\ 5]{frizhairer2020}. The first identity is the one to appreciate: the left-hand side
is a deterministic limit of Riemann sums, computed path by path with no reference to the filtration,
and the right-hand side is an $L^2(\Om)$ limit that exists only up to null sets. That they agree
means the Itô integral, for this class of integrands, admits a pathwise description after all --- at
the price of carrying $\levW^{\mathrm{It\hat o}}$, an object which is itself defined only through
Itô's construction. Rough path theory does not eliminate probability; it isolates it into a single
step, performed once, after which everything is deterministic.

\begin{remark}[The integral depends on the lift]\label{rem:integraldepends}

We can now make precise the claim of Subsection \ref{subsec:typesofrp}, that geometric rough paths
are those on which ordinary calculus survives. The following quantity measures the failure of
\eqref{eq:weaklygeometric}, and turns out to be exactly the correction term in the chain rule.

\end{remark}

\begin{definition}[Bracket of a rough path]\label{def:bracket}
For $\rpX = (X,\lev) \in \RP{\alpha}$ define
\begin{equation}\label{eq:bracket}
[\rpX]_t \;:=\; X_{0,t}\otimes X_{0,t} \;-\; 2\,\Sym\lev_{0,t} \;\in\; V\otimes V .
\end{equation}
\end{definition}

\begin{lemma}\label{lem:bracketincrement}
The bracket is a genuine path, in the sense that its increments are given by the corresponding
two-parameter expression:
\begin{equation}\label{eq:bracketincrement}
[\rpX]_t - [\rpX]_s \;=\; X_{s,t}\otimes X_{s,t} - 2\,\Sym\lev_{s,t} ,
\end{equation}
and $[\rpX] \in \Holder{2\alpha}$.
\end{lemma}

\textbf{\textit{Proof:}} Expanding $X_{0,t} = X_{0,s}+X_{s,t}$ in the first term of
\eqref{eq:bracket},
\[
X_{0,t}\otimes X_{0,t} - X_{0,s}\otimes X_{0,s}
= X_{0,s}\otimes X_{s,t} + X_{s,t}\otimes X_{0,s} + X_{s,t}\otimes X_{s,t},
\]
while Chen's relation \eqref{eq:chenaxiom} gives
$\lev_{0,t}-\lev_{0,s} = \lev_{s,t} + X_{0,s}\otimes X_{s,t}$, so that
\[
2\Sym\lev_{0,t} - 2\Sym\lev_{0,s} = 2\Sym\lev_{s,t} + X_{0,s}\otimes X_{s,t} + X_{s,t}\otimes X_{0,s}.
\]
Subtracting, the two cross terms cancel and \eqref{eq:bracketincrement} follows. The regularity is
then immediate from Definition \ref{def:roughpath}, both terms on the right of
\eqref{eq:bracketincrement} being of order $|t-s|^{2\alpha}$. \qedfilledtext

\medskip

By construction $\rpX$ is weakly geometric precisely when $[\rpX] \equiv 0$; Definition
\ref{def:geometric} could equally well have been phrased that way.

\begin{proposition}[Rough chain rule]\label{prop:roughchainrule}
Let $\rpX = (X,\lev) \in \RP{\alpha}$ with $\alpha\in(\tfrac13,\tfrac12]$, and let
$F \in \mathcal{C}^3_b(V,W)$. Then $\bigl(F(X), DF(X)\bigr) \in \Ctrl{X}$, the rough integral
$\int_0^t DF(X_r)\,\dd\rpX_r$ is defined by Theorem \ref{thm:gubinelli}, and
\begin{equation}\label{eq:roughchainrule}
F(X_t) - F(X_0) \;=\; \int_0^t DF(X_r)\,\dd\rpX_r \;+\; \frac12\int_0^t D^2F(X_r)\,\dd[\rpX]_r ,
\end{equation}
the second integral being a Young integral, well defined since $D^2F(X)\in\Holder{\alpha}$,
$[\rpX]\in\Holder{2\alpha}$ and $3\alpha>1$.

In particular, if $\rpX$ is weakly geometric then $[\rpX]\equiv0$ and \eqref{eq:roughchainrule}
reduces to the classical chain rule
\begin{equation}\label{eq:classicalchainrule}
F(X_t) - F(X_0) \;=\; \int_0^t DF(X_r)\,\dd\rpX_r .
\end{equation}
\qedopen
\end{proposition}

A proof is in \cite[Ch.\ 7]{frizhairer2020}; the mechanism is a second-order Taylor expansion of
$F$ along the partition, in which the symmetric part of $\lev$ is exactly what the quadratic term
of the expansion produces, so that the two agree and cancel precisely when
\eqref{eq:weaklygeometric} holds.

The corollary is the one we have been aiming at since Subsection \ref{subsec:typesofrp}.

\begin{corollary}[It\^o's formula]\label{cor:itoformula}
Let $W$ be a $d$-dimensional Brownian motion and $F\in\mathcal{C}^3_b$. Applying
\eqref{eq:roughchainrule} to the two lifts of Example \ref{ex:brownianroughpath}:
\begin{itemize}
    \item[i)] for $\mathbf{W}^{\mathrm{Strat}}$, which is geometric, one obtains
    $F(W_t)-F(W_0) = \int_0^t DF(W_r)\circ\dd W_r$, the classical chain rule of Stratonovich
    calculus;
    \item[ii)] for $\mathbf{W}^{\mathrm{It\hat o}}$ one computes from \eqref{eq:bracketincrement}
    and the identity $2\Sym\levW^{\mathrm{It\hat o}}_{s,t} = W_{s,t}\otimes W_{s,t} - (t-s)\mathrm{Id}$
    established in Example \ref{ex:brownianroughpath} that
    $[\mathbf{W}^{\mathrm{It\hat o}}]_t = t\,\mathrm{Id}$, whence
    \begin{equation}\label{eq:itofromrough}
    F(W_t)-F(W_0) = \int_0^t DF(W_r)\,\dd W_r + \frac12\int_0^t \Delta F(W_r)\,\dd r ,
    \end{equation}
    which is the case $V(t,x) = F(x)$, $X = W$ of Theorem \ref{thm:itomulti} of Chapter
    \ref{ch:foundations}.
\end{itemize}
\end{corollary}

The qualification matters, and the accounting is worth making explicit. Theorem
\ref{thm:itomulti} covers $\theta(t,X_t)$ for a general $n$-dimensional Itô process and asks only
that $\theta \in \mathcal{C}^{1,2}$; what we have recovered is the time-homogeneous case, driven by
Brownian motion, under the stronger hypothesis $F \in \mathcal{C}^3_b$. The rough-path route is
therefore not a strengthening of Itô's formula but a \emph{re-derivation of a special case of it on
different foundations}, and the extra derivative is the price of those foundations, exactly as in
Remark \ref{rem:c3}. What is bought in exchange is that \eqref{eq:itofromrough} now holds pathwise,
for a fixed trajectory and its lift, with both sides continuous in the rough path metric --- neither
of which the classical statement provides.

\medskip

It is worth being explicit about what has just been shown, because it is the cleanest summary of
the chapter's thesis. In Chapter \ref{ch:foundations}, the second-order term of It\^o's formula arose from
the non-vanishing quadratic variation of Brownian motion, and appeared as an irreducible feature of
stochastic calculus --- the price of integrating against a martingale. Here it has arisen from the
bracket \eqref{eq:bracket} of a particular \emph{lift}, and the Stratonovich lift of the same
trajectory has no such term at all. The correction was never a property of Brownian motion. It was
a property of a choice, made implicitly in Subsection \ref{subsec:itointegral} when the integrand
was evaluated at the left endpoint of each subinterval, and made explicitly here by writing down
$\levW^{\mathrm{It\hat o}}$ rather than $\levW^{\mathrm{Strat}}$.

Let $\rpX = (X,\lev)$ and $\widetilde\rpX = (X, \lev + \increment{f})$ be two lifts of the same
path, as in Proposition \ref{prop:liftnotunique}. Directly from \eqref{eq:roughintegral},
\begin{equation}\label{eq:integraldifference}
\int_s^t Y\,\dd\widetilde\rpX = \int_s^t Y\,\dd\rpX + \int_s^t \Gub_r\,\dd f_r ,
\end{equation}
the last integral existing in the Young sense since $\Gub \in \Holder{\alpha}$,
$f\in\Holder{2\alpha}$ and $3\alpha>1$. Taking $f_t = \tfrac12 t\,\mathrm{Id}$ and $\Gub = DF(W)$
recovers, by Example \ref{ex:brownianroughpath}, the classical relation
$\int F(W)\circ\dd W = \int F(W)\,\dd W + \tfrac12\int F'(W)\,\dd r$ of Chapter \ref{ch:foundations}. The
Itô--Stratonovich correction is thus not a feature of stochastic calculus at all: it is the general
formula \eqref{eq:integraldifference} evaluated at a particular $f$.


\subsection{Differential Equations}\label{subsec:rdes}

With an integral in hand, differential equations are defined in the only way available: as integral
equations. What requires care is that the unknown is now a \emph{pair}.

\begin{definition}[Rough differential equation]\label{def:rde}
Let $\rpX \in \RP{\alpha}$, let $\xi \in W$ and let $f : W \to \Lloc{V,W}$. A pair
$(Y,\Gub) \in \Ctrl{X}$ is a \textbf{solution} of the rough differential equation
\begin{equation}\label{eq:rde}
\dd Y_t = f(Y_t)\,\dd\rpX_t, \qquad Y_0 = \xi,
\end{equation}
if $\Gub = f(Y)$ and, for every $t \in [0,\Mat]$,
\begin{equation}\label{eq:rdeintegral}
Y_t = \xi + \int_0^t f(Y_r)\,\dd\rpX_r ,
\end{equation}
the integral being that of Theorem \ref{thm:gubinelli}, which is well defined because
$\bigl(f(Y), Df(Y)\Gub\bigr) \in \Ctrl{X}$ by Example \ref{ex:controlledexamples} ii).
\end{definition}

\begin{remark}[Why $\Gub = f(Y)$ is imposed]\label{rem:whygubisf}
The constraint is not cosmetic; without it uniqueness fails outright. Consider $X\in\Holder{1}$
lifted canonically, so that $|X_{s,t}| \lesssim |t-s|$. Then for \emph{any} choice of $\Gub$ the
term $\Gub_s X_{s,t}$ is $O(|t-s|)$ and can be absorbed into a $2\alpha$-remainder, so every $\Gub$
makes $(Y,\Gub)$ controlled and \eqref{eq:rdeintegral} carries no information about it. Requiring
$\Gub = f(Y)$ pins the derivative to the equation and restores well-posedness. The reader may
compare Definition \ref{def:sdesolution}, where the analogous role is played by the requirement
that the solution be adapted.
\end{remark}

There is an equivalent formulation, due to Davie, which dispenses with controlled paths entirely
and which is the one a numerical analyst would write down.

\begin{proposition}[Davie's formulation]\label{prop:davie}
A path $Y \in \Holder{\alpha}$ solves \eqref{eq:rde} in the sense of Definition \ref{def:rde} if and
only if
\begin{equation}\label{eq:davie}
Y_{s,t} \;=\; f(Y_s)\,X_{s,t} \;+\; Df(Y_s)f(Y_s)\,\lev_{s,t} \;+\; o\bigl(|t-s|\bigr)
\end{equation}
uniformly in $s\leq t$. \qedopentext
\end{proposition}

The proof is in \cite[Prop.\ 8.8]{frizhairer2020}. Expression \eqref{eq:davie} is worth recognising:
truncated after the first term it is the Euler scheme, and as it stands it is the \emph{Milstein
scheme}, with $\lev_{s,t}$ playing the role of the iterated integral that the Milstein correction
requires. Rough path theory therefore does not merely license the numerical schemes used in
practice; it identifies the object those schemes were always approximating, and Theorem
\ref{thm:ult} will supply their convergence. We shall use this in Chapter \ref{ch:applications}.

\subsection{Existence}\label{subsec:rdeexistence}

\begin{theorem}[Local existence and uniqueness]\label{thm:rdeexistence}
Let $\alpha\in(\tfrac13,\tfrac12]$, let $\rpX \in \RP{\alpha}([0,\Mat],V)$ and let
$f \in \mathcal{C}^3_b\bigl(W,\Lloc{V,W}\bigr)$. Then for every $\xi \in W$ there exist
$\Mat_0 \in (0,\Mat]$ and a unique $(Y,\Gub) \in \Ctrl{X}$ on $[0,\Mat_0]$ solving \eqref{eq:rde}.
The time $\Mat_0$ depends only on $\alpha$, $\lVert f\rVert_{\mathcal{C}^3_b}$ and
$\vertiii{\rpX}_\alpha$, and \emph{not} on $\xi$; consequently the solution extends to a unique
global solution on $[0,\Mat]$.
\end{theorem}

The proof is given in Appendix \ref{app:proof:rdeexistence}. Because the mechanism matters more
than the estimates, we describe its shape here.

One defines the map
\begin{equation}\label{eq:picardmap}
\mathcal{M}(Y,\Gub) \;:=\; \Bigl( \xi + \int_0^{\cdot} f(Y_r)\,\dd\rpX_r,\ f(Y) \Bigr)
\end{equation}
on the Banach space $\Ctrl{X}$, whose fixed points are by Definition \ref{def:rde} exactly the
solutions. One then shows that $\mathcal{M}$ maps a suitable closed ball into itself and is a
contraction on it, provided the time horizon is small; the Banach fixed point theorem does the rest.
This is precisely the Picard iteration by which Theorem \ref{thm:sdeexistence} was proved in Chapter
\ref{ch:foundations}, and the analogy is exact --- but the argument is now entirely deterministic, and
takes place in the space of controlled paths rather than in $L^2(\Om)$.

The contraction factor comes from the horizon. Every estimate in Theorem \ref{thm:gubinelli} carries
a factor $\Mat_0^{\,\alpha}$ when restricted to $[0,\Mat_0]$, and choosing $\Mat_0$ small makes that
factor as small as desired. The subtle point --- and the reason the theorem is stated as it is --- is
that the admissible $\Mat_0$ must be \emph{independent of the initial condition}, since otherwise the
successive restarts used to reach $\Mat$ might have horizons shrinking to zero and the solution would
fail to be global. That independence is secured by the fact, visible in \eqref{eq:roughintegralestimate},
that the estimates involve $\lVert \Gub\rVert_\alpha$ and $\lVert R^Y\rVert_{2\alpha}$ but never
$|Y_0|$, together with the boundedness of $f$. With $\Mat_0$ uniform, one solves on $[0,\Mat_0]$,
restarts at $Y_{\Mat_0}$, and repeats finitely often.

\begin{remark}[On the regularity of $f$]\label{rem:c3}
Three bounded derivatives are required, one more than the two that Example
\ref{ex:controlledexamples} ii) needs to make $f(Y)$ controlled. The extra derivative is spent on
the contraction estimate, where one must compare $f(Y)$ with $f(\widetilde Y)$ \emph{as controlled
paths}, which involves $Df$ and its Lipschitz constant. This is the analogue of the Lipschitz
hypothesis in Definition \ref{def:lipschitzgrowth}; that a rougher driver should demand a smoother
coefficient is the recurring trade of this subject. Without boundedness one still obtains a maximal
solution, which may explode in finite time exactly as in the ODE case.
\end{remark}

\subsection{Uniqueness and Continuity}\label{subsec:ult}

We arrive at the theorem for which the whole apparatus was built. It answers the question raised in
Subsection \ref{subsec:whatfails}, and it answers it affirmatively.

\begin{theorem}[Universal Limit Theorem; continuity of the Itô--Lyons map]\label{thm:ult}
Let $\alpha \in (\tfrac13,\tfrac12]$ and $f \in \mathcal{C}^3_b$. For $\rpX \in \RP{\alpha}$ and
$\xi \in W$ let $\Phi(\xi,\rpX) := Y$ denote the unique solution of \eqref{eq:rde} given by Theorem
\ref{thm:rdeexistence}. Then $\Phi$ is \emph{locally Lipschitz continuous}: for every $M<\infty$
there is $C = C(M,\alpha,f)$ such that, whenever
$\vertiii{\rpX}_\alpha, \vertiii{\widetilde\rpX}_\alpha \leq M$,
\begin{equation}\label{eq:ult}
\bigl\lVert \Phi(\xi,\rpX) - \Phi(\widetilde\xi,\widetilde\rpX) \bigr\rVert_{\alpha}
\;\leq\; C\Bigl( \bigl|\xi - \widetilde\xi\bigr| \;+\; \rpdist\bigl(\rpX,\widetilde\rpX\bigr) \Bigr).
\end{equation}
\end{theorem}

The proof is given in Appendix \ref{app:proof:ult}; it combines the stability of rough integration
(Theorem \ref{thm:integralstability}) with the stability of composition, applied to the fixed point
characterisation \eqref{eq:picardmap}.

\medskip

The significance of \eqref{eq:ult} is best measured against what Part \ref{part:part1} could say.
There, the solution of an SDE was produced as an equivalence class of random variables, obtained
through an $L^2(\Om)$ limit; the map from the driving Brownian path to the solution path existed
only as a measurable map defined up to null sets, and Example \ref{ex:lyonscounterexample} showed
that it could not be continuous in the uniform topology. Theorem \ref{thm:ult} says that once the
driving signal is \emph{enhanced} --- once we carry $\lev$ alongside $X$ and measure distance with
$\rpdist$ --- the solution map becomes not merely continuous but locally Lipschitz, and entirely
deterministic. All the probability in the theory of stochastic differential equations has been
compressed into the single measurable step of constructing $\levW$, performed once in Example
\ref{ex:brownianroughpath}. Everything downstream of it is analysis.

Three consequences deserve mention.

\begin{corollary}[Wong--Zakai]\label{cor:wongzakai}
Let $W$ be a Brownian motion, $W^{(n)}$ its piecewise-linear interpolation over a partition of mesh
tending to zero, and let $Y^{(n)}$ solve the ordinary differential equation
$\dd Y^{(n)} = f(Y^{(n)})\,\dd W^{(n)}$. Then $Y^{(n)}$ converges, almost surely and in
$\alpha$-Hölder norm, to the solution of the \emph{Stratonovich} equation
$\dd Y = f(Y)\circ \dd W$.
\end{corollary}

\textbf{\textit{Proof:}} The canonical lifts of $W^{(n)}$ converge in $\rpdist$ to
$\mathbf{W}^{\mathrm{Strat}}$ --- this is the content of \cite[Ch.\ 3]{frizhairer2020}, and the
Stratonovich rather than the Itô lift appears because each $W^{(n)}$ is smooth, hence its lift is
geometric, hence so is the limit by Proposition \ref{prop:geometricinclusion}. Theorem
\ref{thm:ult} then transfers the convergence to the solutions, and Proposition
\ref{prop:consistency} identifies the limit. \qedfilledtext

\medskip

This is a genuinely non-trivial classical theorem, and here it is a two-line corollary. It also
explains the phenomenon that made it famous: approximating Brownian motion by smooth paths and
solving the resulting ODEs does \emph{not} converge to the Itô solution, because the lifts converge
to the geometric one. In the language of this chapter, the limit remembers the area, and the area of
a smooth approximation is the Stratonovich one.

Second, Theorem \ref{thm:ult} makes numerical analysis possible. The Milstein scheme
\eqref{eq:davie} converges because the map is continuous, and the rate is controlled by
\eqref{eq:ult}; no such argument is available within Itô theory, where discretisation error must be
estimated afresh in $L^2(\Om)$ for each scheme.

Third --- and this is the reason the theory is called \emph{universal} --- nothing in Theorem
\ref{thm:ult} refers to Brownian motion, to a filtration, or to a probability measure. Any process
that can be lifted to $\RP{\alpha}$ inherits the entire theory at once: fractional Brownian motion
for $\Hurst>\tfrac13$ by Theorem \ref{thm:fbmlift}, Gaussian processes with sufficiently regular
covariance, Markov processes, and the piecewise-linear paths formed from observed data. The
existence, uniqueness and continuity results of this section were obtained \emph{once}, and they
apply to all of them.

\section{The Signature Transform}\label{sec:signature}

The rough paths of Definition \ref{def:roughpath} carry two levels because two are what an
$\alpha \in (\tfrac13,\tfrac12]$ \emph{integrator} requires --- the regularity of the integrand is
governed instead by the controlled-path condition \eqref{eq:controlleddef}. Nothing prevents us from
carrying more, and if we carry all of them we obtain an object of independent interest, and the
natural direction in which the present work extends.

\begin{theorem}[Lyons' extension theorem]\label{thm:lyonsextension}
Let $\rpX = (X,\lev) \in \RP{\alpha}$ with $\alpha \in (\tfrac13,\tfrac12]$. Then there is a unique
extension of $\rpX$ to a family
$\bigl(1, X^1_{s,t}, X^2_{s,t}, X^3_{s,t},\dots\bigr)$, with $X^1 = X$, $X^2 = \lev$ and
$X^n_{s,t} \in V^{\otimes n}$, which satisfies Chen's relation at every level,
\begin{equation}\label{eq:chenalllevels}
X^{n}_{s,t} = \sum_{k=0}^{n} X^{k}_{s,u}\otimes X^{n-k}_{u,t}, \qquad s\leq u\leq t,
\end{equation}
and the bounds $|X^n_{s,t}| \leq C_n\,|t-s|^{n\alpha}$ for every $n$. \qedopentext
\end{theorem}

The theorem is Lyons' \cite{lyons1998}; see also \cite[Ch.\ 3]{geng2021}. Its content is that the
higher levels are \emph{not} additional data: once two are given, the rest are determined, and
determined by a construction which is again an instance of the sewing lemma. This is the precise
sense in which Definition \ref{def:roughpath} loses nothing by truncating.

\begin{definition}[Signature]\label{def:signature}
The \textbf{signature} of a rough path $\rpX$ over $[0,\Mat]$ is the full collection of its iterated
integrals over the whole interval,
\begin{equation}\label{eq:signature}
\Sig(\rpX) \;:=\; \bigl(1,\ X^1_{0,\Mat},\ X^2_{0,\Mat},\ X^3_{0,\Mat},\ \dots\bigr),
\end{equation}
and its \textbf{truncation at level $N$}, written $\Sig^N(\rpX)$, is the finite collection
$\bigl(1, X^1_{0,\Mat},\dots,X^N_{0,\Mat}\bigr) \in \Tens{N}(V)$. For a path $x$ of bounded
variation these are the ordinary iterated integrals
\begin{equation}\label{eq:signaturesmooth}
X^n_{0,\Mat} = \int_{0<t_1<\cdots<t_n<\Mat} \dd x_{t_1}\otimes\cdots\otimes \dd x_{t_n}.
\end{equation}
\end{definition}

\begin{example}\label{ex:signatureonedim}
For $d=1$ the tensor products are ordinary products and \eqref{eq:signaturesmooth} evaluates to
$X^n_{0,\Mat} = (x_\Mat - x_0)^n/n!$, independently of the trajectory. The signature of a scalar
path is therefore $\exp(x_\Mat - x_0)$ and carries no information beyond the total increment. All
the content of the transform lies in the non-commutativity of $\otimes$, which is to say in
dimension $d\geq2$.
\end{example}

Three properties make the signature useful, and they are worth stating separately because each
answers a different practical objection.

\begin{proposition}[Factorial decay]\label{prop:factorialdecay}
For a path $x$ of bounded variation with total variation $L := V_1(x)_\Mat$,
\begin{equation}\label{eq:factorialdecay}
\bigl| X^n_{0,\Mat} \bigr| \;\leq\; \frac{L^n}{n!} .
\end{equation}
\end{proposition}

\textbf{\textit{Proof:}} Bounding \eqref{eq:signaturesmooth} termwise, the integrand has modulus at
most $|\dd x_{t_1}|\cdots|\dd x_{t_n}|$, and the domain of integration is the simplex
$0<t_1<\cdots<t_n<\Mat$, which is the $1/n!$ fraction of the cube on which the unrestricted product
integrates to $L^n$. \qedfilledtext

\medskip

The estimate answers the objection that \eqref{eq:signature} is an infinite object: the terms decay
factorially, so a truncation at modest $N$ --- in practice $N \in \{3,4,5\}$ --- retains almost all
of the information. This is what makes $\Sig^N$ a usable feature set.

\begin{proposition}[Shuffle product and linearisation]\label{prop:shuffle}
Let $\rpX$ be geometric. Write $X^{n;I}_{0,\Mat}$ for the coordinate of $X^{n}_{0,\Mat}$ along the
basis vector $e_{i_1}\otimes\cdots\otimes e_{i_n}$, indexed by the word $I = (i_1,\dots,i_n)$. Then
for any two words $I = (i_1,\dots,i_m)$ and $J = (j_1,\dots,j_n)$ in
the letters $\{1,\dots,d\}$,
\begin{equation}\label{eq:shuffle}
X^{m;I}_{0,\Mat}\cdot X^{n;J}_{0,\Mat} \;=\; \sum_{K \,\in\, I \shuffle J} X^{m+n;K}_{0,\Mat},
\end{equation}
the sum being over all $\binom{m+n}{m}$ interleavings of $I$ and $J$ preserving the internal order
of each. \qedopentext
\end{proposition}

Relation \eqref{eq:shuffle} --- the \emph{shuffle product formula}, proved in \cite[Lem.\ 1.15]{geng2021}
by splitting the product of two simplices into the disjoint union of the simplices indexed by the
shuffles --- has a consequence out of all proportion to its appearance. It says that the product of
two coordinates of the signature is a \emph{linear combination} of other coordinates. Hence every
polynomial in the entries of $\Sig(\rpX)$ is again a linear functional of $\Sig(\rpX)$, and any
continuous functional of the path can be approximated, by Stone--Weierstrass, by a \emph{linear}
functional of its signature. A nonlinear problem on path space becomes a linear problem on signature
space. This is the theoretical basis for the use of signatures as features in statistical learning;
see \cite{chevyrevkormilitzin2016} and \cite{levinlyonsni2013}.

It remains to be shown that nothing is lost.

\begin{definition}[Tree-like paths]\label{def:treelike}
A \emph{real tree} is a metric space in which any two points are joined by a unique
non-self-intersecting path, which is moreover a geodesic. A continuous path $x:[0,\Mat]\to V$ is
\textbf{tree-like} if there exist a real tree $\mathcal{T}$ and continuous maps
$\phi : [0,\Mat]\to\mathcal{T}$ and $\psi:\mathcal{T}\to V$ with $\phi(0)=\phi(\Mat)$ and
$x = \psi\circ\phi$. A path is \textbf{tree-reduced} if it contains no tree-like piece.
\end{definition}

The second notion is the one that makes Theorem \ref{thm:signatureuniqueness} an honest
bijection: the signature is injective not on all paths but on the tree-reduced ones, which are
exactly the representatives of the equivalence classes the theorem describes.

Informally, a tree-like path is one that retraces its own steps and thereby cancels itself: it goes
out along some route and returns along the same route, enclosing no area. Such excursions contribute
nothing to any iterated integral, and the following theorem says that they are the \emph{only} thing
the signature forgets.

\begin{theorem}[Uniqueness of the signature]\label{thm:signatureuniqueness}
Let $\rpX$ be a weakly geometric rough path. Then $\Sig(\rpX) = \mathbf{1}$ if and only if
$t\mapsto X_{0,t}$ is tree-like. Consequently --- applying this to
$\rpX \star \widetilde\rpX^{-1}$, the concatenation of $\rpX$ with the time-reversal of
$\widetilde\rpX$, whose signature is $\Sig(\rpX)\otimes\Sig(\widetilde\rpX)^{-1}$ by the group
structure of $T((V))$ --- two weakly geometric rough paths have the same
signature if and only if they coincide up to tree-like pieces. \qedopentext
\end{theorem}

The bounded-variation case is due to Hambly and Lyons \cite{hamblylyons2010}; the extension to rough
paths is due to Boedihardjo, Geng, Lyons and Yang \cite{boedihardjo2016}. A sketch of the argument,
which proceeds by showing that the space of signatures is itself a real tree, is in
\cite[§4.3]{geng2021}.

\medskip

We can now close a loop opened in Chapter \ref{ch:foundations}. There, discussing the log-normal
distribution, we observed that a probability law need not be determined by its moments --- Heyde's
construction, recorded in Appendix \ref{app:heyde}, exhibits distinct densities sharing every
moment with the log-normal --- and remarked that ``knowing every moment of an object is strictly
weaker than knowing the object'', promising that for \emph{paths} the analogous statement could be
made to go the right way. Theorem \ref{thm:signatureuniqueness} is that statement. The iterated
integrals are the moments of a path; unlike the moments of a random variable, they determine it,
and the only ambiguity is the geometrically trivial one of tree-like excursions, which can be
removed outright by the standard device of \emph{time augmentation} --- replacing $x_t$ by
$(t, x_t)$, whose first coordinate is strictly increasing and which therefore admits no tree-like
piece at all. The same device removes the invariance of \eqref{eq:signaturesmooth} under
reparametrisation, which is otherwise a genuine loss of information whenever the speed at which a
path is traversed carries meaning --- as it does, emphatically, for a price series.

The analogy can in fact be pushed one step further, and it is worth doing so, because the promise
made in Chapter \ref{ch:foundations} was about \emph{laws} and Theorem \ref{thm:signatureuniqueness} is
about a fixed path. The statement that genuinely corresponds to ``the moments determine the
distribution'' is obtained by averaging.

\begin{definition}[Expected signature]\label{def:expectedsignature}
Let $\rpX$ be a random geometric rough path, that is, a random variable with values in
$\RPg{\alpha}([0,\Mat],V)$. Its \textbf{expected signature} is the element of the full tensor
algebra $T((V)) = \prod_{n\geq0}V^{\otimes n}$ of Definition \ref{def:tensoralgebra}, obtained by
averaging the signature level by level,
\begin{equation}\label{eq:expectedsignature}
\E\bigl[\Sig(\rpX)\bigr] \;=\; \Bigl(1,\ \E\bigl[X^1_{0,\Mat}\bigr],\ \E\bigl[X^2_{0,\Mat}\bigr],\ \dots\Bigr),
\end{equation}
whenever each expectation exists.
\end{definition}

\begin{theorem}[Chevyrev--Lyons]\label{thm:expectedsignature}
Let $\rpX$ be a random geometric rough path whose expected signature \eqref{eq:expectedsignature}
has infinite radius of convergence, in the sense that
$\sum_{n\geq0}\lambda^{n}\,\bigl|\E[X^n_{0,\Mat}]\bigr| < \infty$ for every $\lambda>0$. Then
$\E[\Sig(\rpX)]$ determines the law of $\rpX$, up to the tree-like equivalence of Theorem
\ref{thm:signatureuniqueness}. \qedopentext
\end{theorem}

The result is due to Chevyrev and Lyons \cite{chevyrevlyons2016}. It is the precise analogue, for
measures on path space, of the elementary fact that a compactly supported law on $\R$ is determined
by its moment sequence --- and the growth hypothesis plays exactly the role that Carleman's condition
played in Chapter \ref{ch:foundations}. The parallel is worth stating in full, since it closes the loop
opened there on the nose rather than approximately:

At the level of individual objects the correspondence is between the power $X^k$ of a random
variable and the iterated integral $X^n_{0,\Mat}$ of a path; there it is Theorem
\ref{thm:signatureuniqueness} that has no counterpart on the left, since the powers of a single
real number determine it trivially. The table below is therefore stated one level up, at the level
of \emph{laws}, which is where the two theories are genuinely parallel.

\begin{center}
\begin{tabular}{@{}p{0.44\linewidth}p{0.44\linewidth}@{}}
\hline
\textbf{Random variables (Ch.\ \ref{ch:foundations})} & \textbf{Paths (this chapter)} \\
\hline
moments $\E[X^k]$ & expected iterated integrals $\E\bigl[X^n_{0,\Mat}\bigr]$ \\
moment generating function & expected signature $\E[\Sig(\rpX)]$ \\
Carleman's condition $\sum_k m_{2k}^{-1/2k}=\infty$ & infinite radius of convergence of $\E[\Sig(\rpX)]$ \\
\emph{fails} for the log-normal (App.\ \ref{app:heyde}) & holds for a wide class, including Brownian motion on a bounded interval \\
\hline
\end{tabular}
\end{center}

The right-hand column is deliberately headed by $\E[\Sig(\rpX)]$ and not by $\Sig(\rpX)$: the
signature is a property of a single path, whereas the moment generating function is a property of a
law, and it is precisely the averaging of Definition \ref{def:expectedsignature} that puts the two
on the same footing. That is the distinction Theorem \ref{thm:expectedsignature} exists to make.

One word on the direction of the two criteria, since the parallel is exact only when it is stated
correctly. Carleman's condition is \emph{sufficient} for determinacy, not necessary; the log-normal
moments grow as $e^{k^2\sigma^2/2}$ and therefore violate it, but that violation establishes only
that the criterion is silent. What establishes indeterminacy is Heyde's explicit construction
(Appendix \ref{app:heyde}), and it is that construction, not the growth rate, which does the work in
Chapter \ref{ch:foundations}. On the right-hand side the corresponding criterion --- infinite radius
of convergence --- is likewise sufficient, and here it is \emph{satisfied}: the signature decays
factorially by Proposition \ref{prop:factorialdecay}, so the series converges for every $\lambda$
whenever the path length has enough integrability, and Theorem \ref{thm:expectedsignature} applies.
The two situations are thus parallel in form and opposite in outcome, and that is the sense in which,
for paths, the analogous statement goes the right way.

\begin{remark}[Why this matters for calibration]\label{rem:signaturecalibration}
The three results above are what make the signature usable as a feature map for path-dependent
problems. Theorem \ref{thm:signatureuniqueness} guarantees that no information is discarded by the
transform; Proposition \ref{prop:factorialdecay} guarantees that truncation is cheap; Proposition
\ref{prop:shuffle} guarantees that any continuous functional of the path is, to arbitrary accuracy,
a \emph{linear} functional of the truncated signature, so that a learning problem on path space
becomes a regression. And, as noted in Remark \ref{rem:H01outside}, none of the three requires a
lifting theorem: an observed data stream, joined by line segments, is a bounded-variation path, and
\eqref{eq:signaturesmooth} applies to it directly, whatever the regularity of the process it
samples. We do not pursue signature-based calibration in this dissertation --- Chapter
\ref{ch:applications} takes a different route --- but we record the transform here because it is the point
at which rough path theory becomes a practical tool rather than a structural one, and because it is
the natural direction in which the present work extends.
\end{remark}



\section{Rough Volatility Models}\label{sec:roughvolmodels}

Subsection \ref{subsec:roughvol} recorded the empirical finding and wrote down the RFSV model
\eqref{eq:rfsv}, which is the natural time-series description of a rough log-volatility. For
pricing one wants something else: a model specified directly under the pricing measure, rich enough
to fit an implied volatility surface, and tractable enough to be calibrated. The model that meets
these requirements is the one that pays the debt incurred in Subsection \ref{subsec:stochvol}, where
the affine structure of the Heston model was described as ``a structural property rather than a
happy accident'' and its rough generalisation promised.

\begin{definition}[Rough Heston]\label{def:roughheston}
Let $\Hurst\in(0,\tfrac12)$ and set $\alpha_{\Hurst} := \Hurst + \tfrac12 \in (\tfrac12,1)$. The
\textbf{rough Heston model} is
\begin{equation}\label{eq:roughheston}
\begin{cases}
\dd \Spot_t = (\rf-\divy)\,\Spot_t\,\dd t + \sqrt{v_t}\,\Spot_t\,\dd \Wt^{\Spot}_t, \\[4pt]
\displaystyle v_t = v_0 + \frac{1}{\Gamma(\alpha_\Hurst)}\int_0^t (t-s)^{\alpha_\Hurst-1}
\Bigl[ \kappa(\theta - v_s)\,\dd s + \xi\sqrt{v_s}\,\dd \Wt^{v}_s \Bigr],
\end{cases}
\end{equation}
with $\dd\langle \Wt^{\Spot},\Wt^{v}\rangle_t = \rho\,\dd t$ in the sense of Definition
\ref{def:correlatedbm}.
\end{definition}

Setting $\alpha_\Hurst = 1$, that is $\Hurst = \tfrac12$, collapses the convolution kernel to a
constant and recovers the classical Heston system \eqref{eq:heston} exactly. For
$\alpha_\Hurst<1$ the kernel $(t-s)^{\alpha_\Hurst-1}$ is singular at $s=t$, and it is this
singularity --- a fractional integral of order $\alpha_\Hurst$ --- that roughens the variance
path to Hölder regularity $\Hurst^{-}$, in accordance with Proposition \ref{prop:fbmproperties} c).

Two structural consequences follow immediately, and they pull in opposite directions.

The first is a loss. The convolution in \eqref{eq:roughheston} makes $v_t$ depend on the entire
history of $\Wt^v$ and not merely on $v_s$, so \textbf{the variance process is no longer Markov},
and neither is the pair $(\Spot,v)$. Every tool of Subsection \ref{section:itoDiff} therefore
becomes unavailable at once: there is no infinitesimal generator, no Kolmogorov equation, no
Feynman--Kac representation. This is precisely the failure anticipated at the end of that
subsection, where the semigroup identity $P_{t+s} = P_t\circ P_s$ was described as a restatement of
the Markov property and as ``the first of the three approaches to fail'' under memory. It has now
failed.

The second is a survival, and it is what makes the model usable. The \emph{affine} structure ---
the fact that the drift and the squared diffusion coefficient are affine in $v$ --- is untouched by
the convolution. Affineness was what produced the Riccati system \eqref{eq:hestonriccati} in the
classical case, and it produces its analogue here.

\begin{theorem}[El Euch--Rosenbaum]\label{thm:roughhestoncf}
Let $\Fwd$ denote the forward price and write $\tau = T-t$. Under \eqref{eq:roughheston} the
characteristic function of the centred log-forward is
\begin{equation}\label{eq:roughhestoncf}
\charf(u,\tau) \;=\; \Ep{\QQ}\Bigl[ \exp\Bigl(iu \log\tfrac{\Fwd_T}{\Fwd_t}\Bigr)\,\Big|\,\Ft \Bigr]
\;=\; \exp\Bigl( \kappa\theta \int_0^{\tau}\psi(u,s)\,\dd s \;+\; v_t\, I^{1-\alpha_\Hurst}\psi(u,\tau) \Bigr),
\end{equation}
where $I^{1-\alpha_\Hurst}$ denotes the Riemann--Liouville fractional integral of order
$1-\alpha_\Hurst$ and $\psi(u,\cdot)$ is the unique continuous solution of the \textbf{fractional
Riccati equation}
\begin{equation}\label{eq:fractionalriccati}
D^{\alpha_\Hurst}\psi(u,\tau) \;=\; \frac{\xi^2}{2}\,\psi(u,\tau)^2 \;+\; \bigl(\rho\xi\, i u - \kappa\bigr)\psi(u,\tau) \;-\; \frac{u^2+iu}{2},
\qquad I^{1-\alpha_\Hurst}\psi(u,0) = 0,
\end{equation}
$D^{\alpha_\Hurst}$ being the Riemann--Liouville fractional derivative of order $\alpha_\Hurst$.
\qedopentext
\end{theorem}

The result is due to El Euch and Rosenbaum \cite{eleucrosenbaum2019}. Comparison with Chapter
\ref{ch:finance} is the point of stating it. Equation \eqref{eq:fractionalriccati} is
\eqref{eq:hestonriccati} \emph{term for term}: the same quadratic in $\psi$, the same linear
coefficient $\rho\xi iu - \kappa$, the same inhomogeneity $-\tfrac12(u^2+iu)$. The single
difference is that the ordinary derivative $\partial_\tau$ has been replaced by a fractional
derivative of order $\alpha_\Hurst = \Hurst+\tfrac12$, and at $\Hurst = \tfrac12$ the two coincide.
The promise made in Subsection \ref{subsec:stochvol} --- that the generalisation would be ``a short
step rather than a fresh start'' --- is discharged by this observation, and the Fourier inversion
\eqref{eq:fourierinversion} carries over verbatim, since it depends on the characteristic function
and on nothing else.

\begin{remark}[What is gained and what it costs]\label{rem:roughhestoncost}
Three comments, in decreasing order of comfort.

\emph{The model reproduces the skew.} By construction the variance has regularity $\Hurst$, so the
short-maturity at-the-money skew behaves as $\tau^{\Hurst-1/2}$ by \eqref{eq:skewpowerlaw}. This is
the empirical shape, and no classical stochastic volatility model produces it.

\emph{The equation is non-local, hence expensive.} Because $D^{\alpha_\Hurst}$ is an integral
operator, \eqref{eq:fractionalriccati} cannot be advanced by a local time-stepping scheme: each
step requires the entire computed history of $\psi$, so a discretisation over $n$ steps costs
$O(n^2)$ rather than $O(n)$. The standard solver is the Adams predictor--corrector scheme. Since
calibration evaluates \eqref{eq:roughhestoncf} at hundreds of strikes for every trial parameter
vector, this cost is the practical bottleneck, and it is the main reason rough Heston is harder to
deploy than its classical counterpart.

\emph{It is not a theorem of this chapter.} Theorem \ref{thm:roughhestoncf} is proved by
techniques from affine Volterra processes and fractional calculus, not by rough path theory, and
nothing in Sections \ref{sec:roughpaths} to \ref{sec:signature} is used in its proof. This is not
an inconsistency but a division of labour, and Section \ref{sec:beyond} explains why it is
unavoidable: at $\Hurst\approx0.1$ the lifting theorems are not available, so a route that does not
require them is not merely convenient but necessary. What rough path theory supplies here is the
vocabulary --- the regularity exponent, the $p$-variation scale, the precise sense in which the
Markov property has been traded away --- rather than the machinery.
\end{remark}




\section{A Purely Rough Model}\label{sec:roughexample}

Everything in this chapter has so far been either machinery or, in the case of Definition
\ref{def:roughheston}, a model whose tractability comes from elsewhere: Remark
\ref{rem:roughhestoncost} was explicit that no result of Sections \ref{sec:roughpaths} to
\ref{sec:signature} enters the proof of Theorem \ref{thm:roughhestoncf}. It is worth closing with
the opposite kind of object --- the smallest model that is \emph{purely} rough, containing no
semimartingale anywhere, on which every construction of the chapter can be exercised once and every
answer written down in closed form.

\begin{definition}[Purely rough asset]\label{def:purelyrough}
Fix $\Hurst \in (\tfrac13,\tfrac12]$, $\sigma>0$ and $\Spot_0>0$, and take $\rf = \divy = 0$, so
that discounting is trivial. Let $\BH$ be a scalar fractional Brownian motion, let
$\mathbf{B}^{\Hurst} = (\BH,\lev^{\Hurst}) \in \RPg{\alpha}$ with $\alpha<\Hurst$ be its
Coutin--Qian lift (Theorem \ref{thm:fbmlift}), and write
$\sigma\mathbf{B}^{\Hurst} := (\sigma\BH,\ \sigma^2\lev^{\Hurst})$ for its dilation, which is again a
geometric rough path by the scaling recorded after Definition \ref{def:rpmetric} and which we shall
need in Subsection \ref{subsec:roughexamplesignature}. The \textbf{purely rough asset} is the
solution, in the sense of Definition \ref{def:rde}, of the linear rough differential equation
\begin{equation}\label{eq:purelyrough}
\dd \Spot_t \;=\; \sigma\,\Spot_t\,\dd\mathbf{B}^{\Hurst}_t , \qquad \Spot_0 > 0 .
\end{equation}
\end{definition}

Since the driver is scalar, one ambiguity of the theory is absent here and it is worth noticing
which. By \eqref{eq:weaklygeometric} a geometric lift must satisfy
$\Sym\lev^{\Hurst}_{s,t} = \tfrac12(\BH_{s,t})^2$, and in one dimension $\lev^{\Hurst}$ \emph{is} its
own symmetric part, so the geometric lift of a scalar path is unique: the choice which Proposition
\ref{prop:liftnotunique} identified as the central structural fact of the subject is, for
\eqref{eq:purelyrough}, no choice at all. Everything below is therefore canonical, and the
comparison with the It\^o lift made in Proposition \ref{prop:purelyroughsolution} is a comparison
with a \emph{non}-geometric object.

Two caveats, both honest rather than technical. The vector field $f(y) = \sigma y$ is not
$\mathcal{C}^3_b$, so Theorem \ref{thm:rdeexistence} does not apply as stated; linear equations are
covered by a separate and entirely standard argument, in which the solution is constructed on a
large ball and the ball is then exhausted --- the same localisation performed in the proof of
Theorem \ref{thm:sdeexistence} of Chapter \ref{ch:foundations}, and carried out for rough equations
in \cite[Ch.\ 10]{frizhairer2020}. And the range $\Hurst>\tfrac13$ is forced by Theorem
\ref{thm:fbmlift}, whereas Subsection \ref{subsec:roughvol} reports $\Hurst\approx0.1$, where by
Remark \ref{rem:H01outside} no canonical lift exists at all. Equation \eqref{eq:purelyrough} is
therefore an illustration of the theory and not a proposal for a market model --- and the second
half of this section shows that it could not be one at any value of $\Hurst$ whatever.

\subsection{The solution, and where Black--Scholes comes from}\label{subsec:roughexamplesolution}

\begin{proposition}\label{prop:purelyroughsolution}
The unique solution of \eqref{eq:purelyrough} is
\begin{equation}\label{eq:purelyroughsolution}
\Spot_t \;=\; \Spot_0\,\exp\bigl(\sigma\BH_t\bigr), \qquad t \in [0,\Mat],
\end{equation}
with no second-order correction, and consequently
\begin{equation}\label{eq:purelyroughmean}
\E\bigl[\Spot_t\bigr] \;=\; \Spot_0\,\exp\Bigl(\tfrac12\,\sigma^2 t^{2\Hurst}\Bigr).
\end{equation}
If instead the driver is the It\^o lift $\mathbf{W}^{\mathrm{It\hat o}}$ of Example
\ref{ex:brownianroughpath} --- that is, if $\Hurst = \tfrac12$ and the geometric lift is exchanged
for the non-geometric one lying over the \emph{same} trajectory --- the solution becomes
$\Spot_0\exp\bigl(\sigma W_t - \tfrac12\sigma^2 t\bigr)$.
\end{proposition}

\textbf{\textit{Proof:}} Put $F(x) := \Spot_0 e^{\sigma x}$, so that $F' = \sigma F$ and
$F'' = \sigma^2F$. The lift is geometric, so its bracket \eqref{eq:bracket} vanishes identically and
the rough chain rule of Proposition \ref{prop:roughchainrule} takes its classical form
\eqref{eq:classicalchainrule},
\[
F\bigl(\BH_t\bigr) - F\bigl(\BH_0\bigr) \;=\; \int_0^t \sigma F\bigl(\BH_r\bigr)\,\dd\mathbf{B}^{\Hurst}_r ,
\]
which is exactly the integral equation \eqref{eq:rdeintegral} for \eqref{eq:purelyrough}; uniqueness
is the linear case of Theorem \ref{thm:rdeexistence}. For \eqref{eq:purelyroughmean}, Proposition
\ref{prop:fbmproperties} a) gives $\BH_t \sim \Normal(0,t^{2\Hurst})$, and
$\E[e^{\sigma Z}] = e^{\sigma^2\Var[Z]/2}$ is the Gaussian moment generating function already used
for the log-normal law in Subsection \ref{subsec:FundamentalDistr}.

It is worth seeing the same computation in Davie's formulation \eqref{eq:davie}, since there it
becomes a Taylor expansion and nothing else. With $f(y) = \sigma y$ one has $f'f = \sigma^2 y$, and
the scalar geometric lift is $\lev^{\Hurst}_{s,t} = \tfrac12(\BH_{s,t})^2$ as noted above, so the
right-hand side of \eqref{eq:davie} reads
\[
\sigma\Spot_s\,\BH_{s,t} \;+\; \tfrac12\sigma^2\Spot_s\bigl(\BH_{s,t}\bigr)^2 ,
\]
which is precisely the second-order Taylor expansion of $\Spot_0e^{\sigma\BH_t}$ about $s$; the
neglected term is $O\bigl(|t-s|^{3\alpha}\bigr) = o(|t-s|)$ because $3\alpha>1$. The Milstein
correction of \eqref{eq:davie} is here the second term of an exponential series.

For the last claim, Proposition \ref{prop:consistency} identifies the rough integral against
$\mathbf{W}^{\mathrm{It\hat o}}$ with the It\^o integral, so \eqref{eq:purelyrough} reads
$\dd\Spot_t = \sigma\Spot_t\,\dd W_t$ and Example \ref{example:blackscholessde} supplies the
solution. The instructive route is the other one: the bracket of the It\^o lift is
$[\mathbf{W}^{\mathrm{It\hat o}}]_t = t\,\mathrm{Id}$ by Corollary \ref{cor:itoformula} ii), so
\eqref{eq:roughchainrule} retains its second-order term and gives
$\dd\bigl(\Spot_0e^{\sigma W_t}\bigr) = \sigma\bigl(\Spot_0e^{\sigma W_t}\bigr)\dd W_t
+ \tfrac12\sigma^2\bigl(\Spot_0e^{\sigma W_t}\bigr)\dd t$. The exponential therefore acquires a
drift it was not asked for, and the factor $e^{-\sigma^2t/2}$ exists to remove it.
\qedfilledtext

\medskip

Three things are worth extracting, in increasing order of importance.

\emph{The Black--Scholes correction is a property of the lift, not of the model.} The two solutions
in Proposition \ref{prop:purelyroughsolution} come from the \emph{same} differential equation, and
the whole difference between them is the $-\tfrac12\sigma^2t$ that Chapter \ref{ch:foundations}
obtained in Example \ref{example:blackscholessde} by applying It\^o's rule to $\log\Spot$. That term
has now been exhibited three times: as a consequence of evaluating the integrand at the left endpoint
(Subsection \ref{subsec:itointegral}), as the bracket of a chosen lift (Corollary
\ref{cor:itoformula}), and here as the sole discrepancy between two solutions of one equation. It is
the same fact on each occasion.

\emph{There is no drift, and yet the mean grows.} Equation \eqref{eq:purelyroughmean} has
$\E[\Spot_t]$ increasing like $e^{\sigma^2t^{2\Hurst}/2}$ although \eqref{eq:purelyrough} contains no
$\dd t$ term whatever. There is no contradiction, because the decomposition of a process into drift
and martingale is a \emph{semimartingale} notion and $\Spot$ is not a semimartingale: were it one,
then $\log$ being $\mathcal{C}^2$ on $(0,\infty)$ and $\Spot$ strictly positive, It\^o's formula
would make $\BH = \sigma^{-1}\log(\Spot/\Spot_0)$ a semimartingale too, contradicting Theorem
\ref{thm:fbmnotsemimartingale}.

\emph{And therefore the model cannot be traded.} The semimartingale property is preserved under an
equivalent change of measure, so $\Spot$ fails to be a semimartingale under \emph{every} $\QQ\sim\PP$
and is in particular a local martingale under none: no equivalent local martingale measure exists.
By the First Fundamental Theorem \ref{thm:firstFTAP} the model violates NFLVR, and Rogers
\cite{rogers1997} --- the same paper that supplies Theorem \ref{thm:fbmnotsemimartingale} ---
exhibits the arbitrage explicitly. It is realised by strategies which trade at unboundedly high
frequency, and it disappears once a minimal delay between transactions is imposed
\cite{cheridito2003}, in the same spirit as the admissibility constraint of Definition
\ref{def:admissible}.

\medskip

That last paragraph explains a design decision which, seen from outside, looks arbitrary.
\textbf{Roughness cannot be placed in the price.} A rough price is not a semimartingale, and a
market whose price is not a semimartingale is not arbitrage-free. Roughness must therefore be placed
one level down, in the \emph{volatility}, where it alters the law of the price without disturbing its
semimartingale property. That is exactly the shape of Definition \ref{def:roughheston}: there the
asset is driven by an ordinary Brownian motion and remains a semimartingale, while the fractional
kernel acts only on the variance. The architecture of every rough volatility model in the literature
is a consequence of the three paragraphs above.

\subsection{The expected signature}\label{subsec:roughexamplesignature}

The model is simple enough that the transform of Section \ref{sec:signature} can be computed
outright, and doing so exhibits in a few lines each the two facts that section stated abstractly:
that a scalar signature is blind, and that the growth hypothesis of Theorem
\ref{thm:expectedsignature} is a genuine condition which here happens to hold.

\begin{proposition}[Expected signature of the driver]\label{prop:roughexpectedsig}
Let $\sigma\mathbf{B}^{\Hurst}$ be as in Definition \ref{def:purelyrough} and let $X^n_{0,\Mat}$
denote the level-$n$ term of its signature over $[0,\Mat]$. Then
\begin{equation}\label{eq:roughexpectedsiglevels}
\E\bigl[X^n_{0,\Mat}\bigr] \;=\;
\begin{cases}
\dfrac{\sigma^{2m}\,\Mat^{2m\Hurst}}{2^{m}\,m!}, & n = 2m, \\[10pt]
0, & n \text{ odd},
\end{cases}
\end{equation}
and, summing the levels,
\begin{equation}\label{eq:roughexpectedsig}
\E\bigl[\Sig(\sigma\mathbf{B}^{\Hurst})_{0,\Mat}\bigr]
\;=\; \exp\Bigl(\tfrac12\sigma^2\Mat^{2\Hurst}\Bigr)
\;=\; \frac{\E[\Spot_\Mat]}{\Spot_0} .
\end{equation}
Moreover
$\sum_{n\geq0}\lambda^{n}\bigl|\E[X^n_{0,\Mat}]\bigr| = \exp\bigl(\tfrac12\lambda^2\sigma^2\Mat^{2\Hurst}\bigr)$
is finite for every $\lambda>0$, so the hypothesis of Theorem \ref{thm:expectedsignature} is
satisfied.
\end{proposition}

\textbf{\textit{Proof:}} The path is scalar, so Example \ref{ex:signatureonedim} gives
$X^n_{0,\Mat} = \sigma^n(\BH_\Mat)^n/n!$. Since $\BH_\Mat\sim\Normal(0,\Mat^{2\Hurst})$ the odd
moments vanish, while $\E[(\BH_\Mat)^{2m}] = (2m-1)!!\,\Mat^{2m\Hurst}$; dividing by $(2m)!$ and
using $(2m)! = 2^m\,m!\,(2m-1)!!$ yields \eqref{eq:roughexpectedsiglevels}. Summing over $m$ produces
the exponential series, which \eqref{eq:purelyroughmean} identifies with
$\E[\Spot_\Mat]/\Spot_0$. Replacing $\sigma$ by $\lambda\sigma$ throughout gives the last claim.
\qedfilledtext

\medskip

What one is entitled to conclude from \eqref{eq:roughexpectedsig} is, however, almost nothing, and
the reason is instructive. Example \ref{ex:signatureonedim} showed that the signature of a scalar
path depends on its terminal increment alone; by Theorem \ref{thm:signatureuniqueness} every scalar
path is therefore equivalent, up to tree-like pieces, to the straight line joining its endpoints. The
law which Theorem \ref{thm:expectedsignature} recovers from \eqref{eq:roughexpectedsig} is thus the
law of $\Spot_\Mat$ and nothing more: a single log-normal marginal, carrying no information about the
trajectory that produced it. This is precisely the loss that Section \ref{sec:signature} warned of,
and it is repaired there by the same device --- time augmentation.

\begin{proposition}[Expected signature after time augmentation]\label{prop:roughexpectedsigaug}
Let $\widehat{X}_t := \bigl(t,\ \sigma\BH_t\bigr) \in \R^2$ and let
$\widehat{\rpX} = (\widehat X,\widehat\lev)$ be its geometric lift, which exists because the second
component is lifted by Theorem \ref{thm:fbmlift} while the remaining entries are Young integrals
against a path of bounded variation. Then
\begin{equation}\label{eq:augsiglevelone}
\E\bigl[\widehat X_{0,\Mat}\bigr] \;=\; \bigl(\Mat,\; 0\bigr),
\end{equation}
and, writing $\widehat\lev^{ij}$ for the component of $\widehat\lev$ along $e_i\otimes e_j$,
\begin{equation}\label{eq:augsigleveltwo}
\E\bigl[\widehat\lev_{0,\Mat}\bigr] \;=\; \frac12
\begin{pmatrix}
\Mat^{2} & 0 \\[3pt]
0 & \sigma^2\Mat^{2\Hurst}
\end{pmatrix}.
\end{equation}
\end{proposition}

\textbf{\textit{Proof:}} The first level is the increment
$\widehat X_{0,\Mat} = (\Mat,\ \sigma\BH_\Mat)$, and $\E[\BH_\Mat]=0$. For the second level,
geometricity \eqref{eq:weaklygeometric} determines the symmetric part outright and pathwise,
\[
\Sym\widehat\lev_{0,\Mat} \;=\; \tfrac12\,\widehat X_{0,\Mat}\otimes\widehat X_{0,\Mat}
\;=\; \frac12\begin{pmatrix} \Mat^2 & \sigma\Mat\,\BH_\Mat \\[2pt] \sigma\Mat\,\BH_\Mat & \sigma^2(\BH_\Mat)^2 \end{pmatrix},
\]
so the two diagonal entries are already identified, and taking expectations with
$\E[(\BH_\Mat)^2] = \Mat^{2\Hurst}$ gives the diagonal of \eqref{eq:augsigleveltwo}. For the
off-diagonal, the first coordinate has bounded variation, so
$\widehat\lev^{21}_{0,\Mat} = \sigma\int_0^{\Mat}\BH_r\,\dd r$ is an ordinary Riemann integral;
Fubini applies, since
$\E\int_0^{\Mat}|\BH_r|\,\dd r = \sqrt{2/\pi}\int_0^{\Mat} r^{\Hurst}\,\dd r < \infty$, and yields
$\E\bigl[\widehat\lev^{21}_{0,\Mat}\bigr] = \sigma\int_0^\Mat \E[\BH_r]\,\dd r = 0$. The remaining
entry follows from the symmetric part, since
$\widehat\lev^{12}_{0,\Mat} + \widehat\lev^{21}_{0,\Mat} = \sigma\Mat\,\BH_\Mat$ has expectation
zero. \qedfilledtext

\begin{remark}[Where the Hurst parameter lives]\label{rem:Hlivesatleveltwo}
Compare \eqref{eq:augsiglevelone} with \eqref{eq:augsigleveltwo}. The first level of the expected
signature is $(\Mat,0)$ for \emph{every} value of $\Hurst$: it cannot tell a rough asset from a
smooth one. The parameter appears for the first time at level two, in the single entry
$\tfrac12\sigma^2\Mat^{2\Hurst}$, and it appears there as an exponent of the time horizon.

This is the structural reason behind an estimator that Subsection \ref{subsec:roughvol} presented as
an empirical recipe. The scaling relation $m(q,\Delta)\propto\Delta^{q\Hurst}$ of
\eqref{eq:scalingmoment} identifies $\Hurst$ from a \emph{second} moment of increments when $q=2$;
what \eqref{eq:augsigleveltwo} adds is that this is not a matter of convenience but of necessity,
since no first-order functional of the path carries the exponent at all. Roughness is a level-two
phenomenon, in the same sense in which the second level of a rough path was the genuinely new data of
Proposition \ref{prop:liftnotunique}.

Two checks close the discussion. At $\Hurst=\tfrac12$ the entry reads $\tfrac12\sigma^2\Mat$, which
is the level-two term of Fawcett's formula
$\E\bigl[\Sig(\mathbf{W}^{\mathrm{Strat}})_{0,\Mat}\bigr] = \exp\bigl(\tfrac12\Mat\,\mathrm{Id}\bigr)$
for Brownian motion under the Stratonovich lift \cite{frizvictoir2010}. And beyond level two no
formula of Fawcett's type is known for $\Hurst\neq\tfrac12$: the closed forms of this section stop
exactly where the independence of increments stops, which is a fair measure of what the fractional
case costs.
\end{remark}

The example has now used the chapter once through. The lift came from Theorem \ref{thm:fbmlift};
geometricity removed the second-order term from the chain rule, and its failure for the It\^o lift
put the term back and returned Black--Scholes; the equation was solved in the sense of Definition
\ref{def:rde}; the signature of Section \ref{sec:signature} was computed, found blind in one
dimension and informative after time augmentation; and the growth hypothesis of Theorem
\ref{thm:expectedsignature} was verified rather than assumed.

What the example does \emph{not} supply is a model one could price against, and it is worth being
careful about how much that concedes. The obstruction is sharp, but it is also narrow: what fails is
the existence of an equivalent local martingale measure, and an equivalent local martingale measure
is an object required by \emph{pricing} and by nothing else. It is tempting to read Theorem
\ref{thm:fbmnotsemimartingale} as a verdict on the usefulness of rough paths in finance. It is not.
It is a statement about \emph{where} in finance they may be used, and the next section takes that
statement literally: it poses a problem in which no pricing measure appears anywhere, drives it with
a price as illegitimate as \eqref{eq:purelyrough}, and solves it in closed form.



\section{What Rough Paths Are For: Optimal Execution}\label{sec:execution}

Section \ref{sec:roughexample} ended on an obstruction, and it is worth being exact about what the
obstruction forbids. The First Fundamental Theorem \ref{thm:firstFTAP} is a statement about
\emph{pricing measures}: it says that a market admits no free lunch precisely when an equivalent
local martingale measure exists. A rough price has none, so a rough price cannot be the underlying
of an arbitrage-free pricing theory. That is the whole of the damage.

It is not, however, the whole of finance. A large part of the subject consists of problems posed
under the physical measure $\PP$, or under no measure at all, in which no expectation is ever taken
against a pricing kernel: the execution of a large order, the attribution of a hedging profit and
loss, the measurement of risk, the estimation of a parameter, the choice of a portfolio. None of
these is a pricing problem, none of them requires an equivalent local martingale measure, and none
of them is touched by Theorem \ref{thm:fbmnotsemimartingale}. Roughness, meanwhile, is an empirical
fact (Subsection \ref{subsec:roughvol}). The natural conclusion is that rough path theory is a
$\PP$-tool applied to a $\QQ$-innocent class of problems, and this section works one such problem
through to a closed-form answer.

\begin{remark}[Where the lift is, and is not, needed]\label{rem:noliftneeded}
One preliminary observation removes the restriction $\Hurst>\tfrac13$ from everything that follows,
and it is the reason the applications below are unaffected by Section \ref{sec:beyond}. In practice
the signature is computed from an \emph{observed}, discretely sampled trajectory, linearly
interpolated between ticks. Such a path has bounded variation, its iterated integrals
\eqref{eq:signaturesmooth} are ordinary Riemann--Stieltjes integrals, and no lift is required at
all. What rough path theory supplies in that setting is not the construction of the object but the
guarantees attached to it: factorial decay (Proposition \ref{prop:factorialdecay}) is what licenses
truncation, the shuffle relation (Proposition \ref{prop:shuffle}) is what linearises the problem,
Theorem \ref{thm:signatureuniqueness} is what says nothing has been discarded, and the Universal
Limit Theorem \ref{thm:ult} is what says the answer is stable as the sampling refines rather than an
artefact of the grid. The continuum model below is stated with a lift because we wish to compute
expectations in closed form; the procedure it justifies does not need one.
\end{remark}

% ---------------------------------------------------------------------------
\subsection{The execution problem}\label{subsec:executionproblem}

A trader must liquidate $Q_0>0$ shares over a horizon $[0,\Mat]$ --- hours, not years. Selling fast
moves the price against them; selling slowly exposes them to the market. This is the
\emph{optimal execution} problem of Almgren and Chriss \cite{almgrenchriss2001}, and it is the
canonical example of a financial problem in which no pricing measure appears: the objective is an
expected cost under $\PP$, and the trader is not valuing a claim but choosing a schedule.

\begin{definition}[Execution problem]\label{def:execution}
Let $\Spot$ denote the \emph{unaffected} price, let $Q_t$ be the inventory remaining at time $t$ and
let $\pi_t = -\dot Q_t$ be the trading rate, so that $Q_t = Q_0-\int_0^t\pi_r\,\dd r$. Selling at
rate $\pi_t$ realises the price $\Spot_t-\eta\,\pi_t$ with $\eta>0$ the \emph{temporary impact}
coefficient, and any inventory $Q_\Mat$ left at the horizon is liquidated in a block at
$\Spot_\Mat-\zeta\,Q_\Mat$, with $\zeta>0$. The \textbf{implementation shortfall} of the schedule
$\pi$ is
\begin{equation}\label{eq:shortfalldef}
  \mathcal{C}(\pi) \;=\; Q_0\Spot_0
  \;-\;\int_0^{\Mat}\pi_t\bigl(\Spot_t-\eta\pi_t\bigr)\dd t
  \;-\;Q_\Mat\bigl(\Spot_\Mat-\zeta Q_\Mat\bigr),
\end{equation}
the arrival value of the position less what the liquidation actually raises.
\end{definition}

Here $\pi$ is a trading strategy in exactly the sense fixed in the table of notation, and the letter
carries no other meaning in this section; the constant $\pi$ appears only inside Gaussian
normalisations, and never with a time subscript. The coefficient $\zeta$ is a cost, and is
unrelated to the scaling exponents $\zeta_q$ of Subsection \ref{subsec:roughvol}, which always carry
their subscript.

\begin{lemma}[Reduction of the shortfall]\label{lem:shortfall}
For any $C^{1}$ inventory path $Q$ with $Q_0$ given,
\begin{equation}\label{eq:shortfallreduced}
  \mathcal{C}(\pi)\;=\;-\int_0^{\Mat}Q_t\,\dd\Spot_t
  \;+\;\eta\int_0^{\Mat}\pi_t^{2}\,\dd t
  \;+\;\zeta\,Q_\Mat^{2}.
\end{equation}
\end{lemma}

\textbf{\textit{Proof:}} Since $Q$ is $C^{1}$ it has bounded variation, so $\int_0^\Mat Q\,\dd\Spot$
is a Young integral --- the criterion of Remark \ref{rem:criticalexponent} reads
$1+\alpha>1$ here, since $Q$ is Lipschitz and $\Spot\in\Holder{\alpha}$ --- and integration by parts
is available:
$\int_0^{\Mat}\dot Q_t\Spot_t\,\dd t = Q_\Mat\Spot_\Mat-Q_0\Spot_0-\int_0^{\Mat}Q_t\,\dd\Spot_t$.
With $\pi=-\dot Q$ this gives
$\int_0^{\Mat}\pi_t\Spot_t\,\dd t = Q_0\Spot_0-Q_\Mat\Spot_\Mat+\int_0^{\Mat}Q_t\,\dd\Spot_t$.
Substituting into \eqref{eq:shortfalldef}, the two occurrences of $Q_0\Spot_0$ cancel and so do the
two of $Q_\Mat\Spot_\Mat$, leaving \eqref{eq:shortfallreduced}. \qedfilledtext

\medskip

The three terms of \eqref{eq:shortfallreduced} are, in order, the market risk carried by the
unliquidated inventory, the impact paid for trading, and the cost of the terminal block. Only the
first involves the price path, and it involves it through an integral against $\dd\Spot$ --- which is
where the regularity of $\Spot$ enters, and where a semimartingale would ordinarily be assumed.

\begin{definition}[Rough execution model]\label{def:roughexecution}
Fix $\Hurst\in(\tfrac13,\tfrac12]$ and $\sigma>0$, let $\BH$ be a scalar fractional Brownian motion
and set
\begin{equation}\label{eq:arithmeticrough}
  \Spot_t \;=\; \Spot_0+\sigma\,\BH_t ,\qquad t\in[0,\Mat].
\end{equation}
The \textbf{time-augmented driver} is $\widehat X_t=(t,\ \sigma\BH_t)\in\R^{2}$, lifted to the
geometric rough path $\widehat\rpX$ of Proposition \ref{prop:roughexpectedsigaug}. We write
$\mathbb{S}^{I}_t$ for the coordinate of $\Sig(\widehat\rpX)_{0,t}$ along the word $I$ in the letters
$\{1,2\}$, indexed as in Proposition \ref{prop:shuffle}, so that $\mathbb{S}^{1}_t=t$ and
$\mathbb{S}^{2}_t=\sigma\BH_t$.
\end{definition}

The price \eqref{eq:arithmeticrough} is arithmetic rather than exponential, as is standard over an
intraday horizon, and it is exactly as illegitimate a \emph{pricing} model as \eqref{eq:purelyrough}:
it is not a semimartingale, so by the argument of Subsection \ref{subsec:roughexamplesolution} it
admits no equivalent local martingale measure. That is the point. Nothing in Definition
\ref{def:execution} asked for one.

% ---------------------------------------------------------------------------
\subsection{Schedules as linear functionals of the signature}\label{subsec:sigstrategies}

We now restrict the schedule to the class in which the problem linearises.

\begin{definition}[Signature schedule]\label{def:sigschedule}
Let $N\in\N$ and let $\ell$ be a linear functional on $\Tens{N}(\R^{2})$, that is, a finite family of
coefficients $(\ell_I)_{|I|\leq N}$ indexed by words. The associated \textbf{signature schedule} is
\begin{equation}\label{eq:sigschedule}
  \pi_t \;=\; \ell\bigl(\Sig^{N}(\widehat\rpX)_{0,t}\bigr)\;=\;\sum_{|I|\leq N}\ell_I\,\mathbb{S}^{I}_t .
\end{equation}
\end{definition}

The class is not as special as it looks. By Proposition \ref{prop:shuffle} and Stone--Weierstrass,
linear functionals of the signature are dense in the continuous functionals of the path; the family
\eqref{eq:sigschedule} therefore approximates \emph{any} continuous non-anticipative schedule, and
Proposition \ref{prop:factorialdecay} bounds the error committed by stopping at level $N$. What the
restriction buys is that the two operations the objective \eqref{eq:shortfallreduced} performs on
$\pi$ --- integrating it, and squaring it --- both stay inside the class. That is the content of the
next two observations, and everything afterwards is bookkeeping.

\begin{lemma}[Integration is concatenation]\label{lem:concatenation}
For any word $I$ and either letter $j\in\{1,2\}$,
$\mathbb{S}^{Ij}_t=\int_0^t\mathbb{S}^{I}_r\,\dd\widehat X^{\,j}_r$. In particular, writing $\ell j$ for the
functional with coefficients $(\ell j)_{Ij}=\ell_I$ and zero on words not ending in $j$,
\begin{equation}\label{eq:concatenation}
  \int_0^t \pi_r\,\dd r \;=\; (\ell 1)\bigl(\mathbb{S}^{N+1}_t\bigr),
  \qquad
  \int_0^t \pi_r\,\dd\Spot_r \;=\; (\ell 2)\bigl(\mathbb{S}^{N+1}_t\bigr).
\end{equation}
\end{lemma}

\textbf{\textit{Proof:}} The first assertion is the definition \eqref{eq:signaturesmooth} of the
iterated integral read one level at a time. The two displays follow by linearity, using
$\dd\widehat X^{1}_r=\dd r$ and $\dd\widehat X^{2}_r=\sigma\,\dd\BH_r=\dd\Spot_r$.
\qedfilledtext

\medskip

Lemma \ref{lem:concatenation} is the reason the time coordinate was adjoined in Proposition
\ref{prop:roughexpectedsigaug}. Adjoining $t$ does more than break the blindness of a scalar
signature recorded in Example \ref{ex:signatureonedim}: it makes \emph{time integration} an
operation on words, namely right-concatenation with the letter $1$. The inventory, which is the time
integral of the schedule, is therefore again a linear functional of the signature --- one level
higher.

\begin{proposition}[The objective is a quadratic form in the signature]\label{prop:quadform}
Let $\pi$ be the signature schedule \eqref{eq:sigschedule} and write
$\mathbb{S}:=\Sig^{N+2}(\widehat\rpX)_{0,\Mat}$ for the truncated signature of the realised
trajectory. Then, \emph{on that trajectory},
\begin{equation}\label{eq:quadform}
  \mathcal{C}(\pi)\;=\;\ell\transp\!\mathfrak{A}(\mathbb{S})\,\ell\;+\;\mathfrak{b}(\mathbb{S})\transp\ell\;+\;c(\mathbb{S}),
\end{equation}
where, writing $e_I$ for the functional dual to the word $I$,
\begin{align}
  \mathfrak{A}_{IJ}(\mathbb{S}) &= \eta\,\bigl((e_I\shuffle e_J)1\bigr)(\mathbb{S})
            \;+\;\zeta\,\bigl((e_I1)\shuffle(e_J1)\bigr)(\mathbb{S}),\label{eq:Adef}\\
  \mathfrak{b}_{I}(\mathbb{S}) &= (e_I12)(\mathbb{S})\;-\;2\zeta Q_0\,(e_I1)(\mathbb{S}),\label{eq:bdef}
\end{align}
and $c(\mathbb{S})=\zeta Q_0^{2}-Q_0\,\mathbb{S}^{2}$ does not involve $\ell$. Both $\mathfrak{A}$ and
$\mathfrak{b}$ are \emph{linear} in $\mathbb{S}$.
\end{proposition}

\begin{corollary}\label{cor:quadformexpectation}
Writing $\overline{\mathbb{S}}:=\E[\mathbb{S}]$ and taking expectations in \eqref{eq:quadform},
\begin{equation}\label{eq:quadformE}
  \E\bigl[\mathcal{C}(\pi)\bigr]\;=\;\ell\transp\!\mathfrak{A}(\overline{\mathbb{S}})\,\ell
  \;+\;\mathfrak{b}(\overline{\mathbb{S}})\transp\ell\;+\;\zeta Q_0^{2},
\end{equation}
the term $Q_0\E[\mathbb{S}^{2}]$ vanishing because $\E[\BH_\Mat]=0$. The matrix
$\mathfrak{A}(\overline{\mathbb{S}})$ is symmetric and positive definite for $\eta,\zeta>0$, so the minimiser is
unique and is
\begin{equation}\label{eq:optimalell}
  \ell^{*}\;=\;-\tfrac12\,\mathfrak{A}(\overline{\mathbb{S}})^{-1}\,\mathfrak{b}(\overline{\mathbb{S}}).
\end{equation}
\end{corollary}

\textbf{\textit{Proof:}} Take the three terms of \eqref{eq:shortfallreduced} in turn. By Lemma
\ref{lem:concatenation}, $Q_t=Q_0-(\ell 1)(\mathbb{S}_t)$, so
\[
  \int_0^{\Mat}Q_t\,\dd\Spot_t = Q_0\,\mathbb{S}^{2}_\Mat-(\ell 12)(\mathbb{S}_\Mat),
\]
a second application of the lemma with $j=2$. For the impact term, Proposition \ref{prop:shuffle}
gives $\pi_t^{2}=(\ell\shuffle\ell)(\mathbb{S}_t)$, whence
$\int_0^{\Mat}\pi_t^{2}\dd t=\bigl((\ell\shuffle\ell)1\bigr)(\mathbb{S}_\Mat)$, again by Lemma
\ref{lem:concatenation}. For the terminal term,
$Q_\Mat^{2}=Q_0^{2}-2Q_0(\ell1)(\mathbb{S}_\Mat)+\bigl((\ell1)\shuffle(\ell1)\bigr)(\mathbb{S}_\Mat)$. Every
resulting expression is a linear functional of $\mathbb{S}^{N+2}_\Mat$ whose coefficients are quadratic or
linear in $\ell$; taking expectations replaces $\mathbb{S}_\Mat$ by $\overline{\mathbb{S}}$ and collects the terms
into \eqref{eq:Adef}--\eqref{eq:bdef}, the contribution $Q_0\E[\mathbb{S}^{2}_\Mat]$ vanishing because
$\E[\BH_\Mat]=0$. Positive definiteness holds because
$\ell\transp \mathfrak{A}\ell=\eta\,\E\int_0^{\Mat}\pi_t^{2}\dd t+\zeta\,\E\bigl[(\ell1)(\mathbb{S}_\Mat)^{2}\bigr]\geq0$,
with equality only if $\pi\equiv0$ almost surely, which forces $\ell=0$ by the uniqueness of Theorem
\ref{thm:signatureuniqueness}. \qedfilledtext

\medskip

The order of these two statements is the point, and it is worth insisting on. Proposition
\ref{prop:quadform} is an \emph{identity}, valid on each trajectory separately: no expectation is
taken, no measure is used, and the shuffle relation on which it rests is itself pathwise. The
expectation \eqref{eq:quadformE} is a corollary of it and not the other way round. Everything the
method needs is therefore available from realised paths, and the probabilistic model enters only if
one chooses to compute $\overline{\mathbb{S}}$ in closed form rather than measure it.

Three further features are worth naming. The problem is finite-dimensional: the unknown is a vector
of $\sum_{n\leq N}2^{n}$ numbers, and \eqref{eq:optimalell} is a linear solve. It is
\emph{model-free in the coefficients}: $\mathfrak{A}$ and $\mathfrak{b}$ depend on the market only through
$\overline{\mathbb{S}}$, so a change of model is a change of one vector of constants and not a change
of method. And it requires the signature two levels beyond the schedule, since the objective
integrates $\pi$ once against $\dd t$ and once against $\dd\Spot$.

\begin{remark}[Calibration is an average]\label{rem:erm}
Because \eqref{eq:quadform} holds path by path, the coefficients need not be computed from a model at
all. Given $M$ historical liquidation windows with realised signatures
$\mathbb{S}^{(1)},\dots,\mathbb{S}^{(M)}$, put
$\widehat{\mathbb{S}}=\tfrac1M\sum_{m}\mathbb{S}^{(m)}$ and
$\widehat\ell^{*}=-\tfrac12\mathfrak{A}(\widehat{\mathbb{S}})^{-1}\mathfrak{b}(\widehat{\mathbb{S}})$. By the
linearity of $\mathfrak{A}$ and $\mathfrak{b}$ this is exactly the empirical risk minimiser over the class
\eqref{eq:sigschedule}, and it is what \cite{kalsi2020} implements.

The whole of ``calibration'' is therefore an \emph{average of signatures} followed by one linear
solve. There is no parametric family, no likelihood, and no measure: no $\Hurst$ is estimated, no
$\PP$ is posited, and the question of whether the market admits arbitrage --- Remark
\ref{rem:isitarbitrage} --- does not arise, because nothing in the procedure presumes it does not.
By Remark \ref{rem:noliftneeded} the signatures are those of sampled, piecewise-linear paths and no
lift is involved. This should be contrasted with the calibration of Chapter
\ref{ch:applications}, which fits a cross-section of option \emph{prices} at one date under $\QQ$
and to which none of the above applies: the two procedures use different data, answer different
questions, and are available in different markets.
\end{remark}

% ---------------------------------------------------------------------------
\subsection{A worked instance, in closed form}\label{subsec:executionclosed}

Take $N=1$, so that the schedule is
\begin{equation}\label{eq:worked}
  \pi_t \;=\; \ell_\emptyset+\ell_2\,\sigma\BH_t ,
\end{equation}
a constant rate plus a term proportional to the price move so far: sell faster when the price has
risen, or slower, according to the sign of $\ell_2$. By Proposition \ref{prop:quadform} the objective
needs the expected signature to level three. The level-one and level-two entries are supplied by
Proposition \ref{prop:roughexpectedsigaug}; the one new entry is the following, and it is where the
whole answer lives.

\begin{proposition}[A level-three entry]\label{prop:levelthree}
With $\widehat\rpX$ as in Definition \ref{def:roughexecution},
\begin{equation}\label{eq:levelthree}
  \E\bigl[\mathbb{S}^{212}_\Mat\bigr]\;=\;\frac{(2\Hurst-1)\,\sigma^{2}\,\Mat^{2\Hurst+1}}{2\,(2\Hurst+1)} ,
\end{equation}
which vanishes if and only if $\Hurst=\tfrac12$, and is strictly negative for $\Hurst<\tfrac12$.
Moreover
\begin{equation}\label{eq:levelthreeaux}
  \E\Bigl[\bigl(\mathbb{S}^{21}_\Mat\bigr)^{2}\Bigr]=\frac{\sigma^{2}\Mat^{2\Hurst+2}}{2\Hurst+2},
  \qquad
  \E\bigl[\mathbb{S}^{221}_\Mat\bigr]=\frac{\sigma^{2}\Mat^{2\Hurst+1}}{2(2\Hurst+1)} .
\end{equation}
\end{proposition}

\textbf{\textit{Proof:}} Write $G_t:=\int_0^t\BH_r\,\dd r$, so that $\mathbb{S}^{21}_\Mat=\sigma G_\Mat$
and, by Lemma \ref{lem:concatenation}, $\mathbb{S}^{212}_\Mat=\sigma^{2}\int_0^{\Mat}G_t\,\dd\BH_t$. Since
$G$ is $C^{1}$ the integral is a Young integral and integration by parts applies:
\[
  \int_0^{\Mat}G_t\,\dd\BH_t \;=\; G_\Mat\BH_\Mat-\int_0^{\Mat}\BH_t\,\dd G_t
  \;=\; G_\Mat\BH_\Mat-\int_0^{\Mat}\bigl(\BH_t\bigr)^{2}\dd t .
\]
By Proposition \ref{prop:fbmproperties} a) the covariance is
$\E[\BH_s\BH_t]=\tfrac12\bigl(s^{2\Hurst}+t^{2\Hurst}-|t-s|^{2\Hurst}\bigr)$, so Fubini gives
\[
  \E\bigl[G_\Mat\BH_\Mat\bigr]=\int_0^{\Mat}\!\E[\BH_s\BH_\Mat]\,\dd s
  =\tfrac12\Bigl(\tfrac{\Mat^{2\Hurst+1}}{2\Hurst+1}+\Mat^{2\Hurst+1}-\tfrac{\Mat^{2\Hurst+1}}{2\Hurst+1}\Bigr)
  =\tfrac12\Mat^{2\Hurst+1},
\]
while $\E\int_0^{\Mat}(\BH_t)^{2}\dd t=\int_0^{\Mat}t^{2\Hurst}\dd t=\Mat^{2\Hurst+1}/(2\Hurst+1)$.
Subtracting and multiplying by $\sigma^{2}$ yields \eqref{eq:levelthree}. For the first identity in
\eqref{eq:levelthreeaux},
\[
  \E\bigl[G_\Mat^{2}\bigr]=\int_0^{\Mat}\!\!\int_0^{\Mat}\!\E[\BH_s\BH_r]\,\dd s\,\dd r
  =\frac{\Mat^{2\Hurst+2}}{2\Hurst+1}-\frac{\Mat^{2\Hurst+2}}{(2\Hurst+1)(2\Hurst+2)}
  =\frac{\Mat^{2\Hurst+2}}{2\Hurst+2},
\]
using $\int_0^{\Mat}\!\int_0^{\Mat}|s-r|^{2\Hurst}\dd s\,\dd r=2\Mat^{2\Hurst+2}/[(2\Hurst+1)(2\Hurst+2)]$.
The last identity is immediate from $\mathbb{S}^{22}_t=\tfrac12(\mathbb{S}^{2}_t)^{2}$, which is geometricity
\eqref{eq:weaklygeometric} in one dimension, followed by Lemma \ref{lem:concatenation}.
\qedfilledtext

\medskip

\begin{remark}[Two consistency checks]\label{rem:levelthreechecks}
At $\Hurst=\tfrac12$ the driver is Brownian, $G$ is adapted and square-integrable, and
$\int_0^\Mat G\,\dd\Wt$ is a true It\^o integral of mean zero; formula \eqref{eq:levelthree} returns
exactly $0$. Independently, the shuffle relation \eqref{eq:shuffle} applied to the words $(2)$ and
$(2,1)$ gives $\mathbb{S}^{2}\mathbb{S}^{21}=2\,\mathbb{S}^{221}+\mathbb{S}^{212}$, and the three
expectations computed above
satisfy it: $\tfrac12\sigma^{2}\Mat^{2\Hurst+1}$ on the left, and on the right
$\sigma^{2}\Mat^{2\Hurst+1}/(2\Hurst+1)+(2\Hurst-1)\sigma^{2}\Mat^{2\Hurst+1}/[2(2\Hurst+1)]$, which
is the same number. The second check is a genuine one: it uses a structural identity of the
signature to verify a computation performed by Fubini, and the two routes share no step.
\end{remark}

The two entries of Proposition \ref{prop:levelthree} differ by the factor $2\Hurst-1$ and by nothing
else, which has a use of its own.

\begin{corollary}[A scale-free estimator of the Hurst parameter]\label{cor:Hestimator}
With the notation of Proposition \ref{prop:levelthree},
\begin{equation}\label{eq:Hratio}
  \frac{\E\bigl[\mathbb{S}^{212}_\Mat\bigr]}{\E\bigl[\mathbb{S}^{221}_\Mat\bigr]}\;=\;2\Hurst-1,
  \qquad\text{that is}\qquad
  \Hurst\;=\;\frac12\left(1+\frac{\E\bigl[\mathbb{S}^{212}_\Mat\bigr]}{\E\bigl[\mathbb{S}^{221}_\Mat\bigr]}\right).
\end{equation}
\end{corollary}

\textbf{\textit{Proof:}} Divide \eqref{eq:levelthree} by the second identity of
\eqref{eq:levelthreeaux}; the denominator is strictly positive. \qedfilledtext

\medskip

Both the scale $\sigma$ and the horizon $\Mat$ cancel, so \eqref{eq:Hratio} identifies $\Hurst$ from
two level-three entries measured over a \emph{single} window, without the regression across
$\Delta$ that the scaling relation \eqref{eq:scalingmoment} of Subsection \ref{subsec:roughvol}
requires. Both entries are Riemann sums on a sampled path --- $\mathbb{S}^{221}$ is
$\tfrac12\int_0^{\Mat}(\text{level})^{2}\dd t$ and $\mathbb{S}^{212}$ is the Young integral of
Proposition \ref{prop:levelthree} --- so by Remark \ref{rem:noliftneeded} no lift is needed to
evaluate either. The ratio maps $\Hurst\in(0,1)$ affinely onto $(-1,1)$ and equals zero exactly at
$\Hurst=\tfrac12$, so it is simultaneously an estimator and a test of the Brownian null.

\begin{remark}[What it should be applied to, and what may go wrong]\label{rem:Hestimatorscope}
Three cautions, and the first is the one that matters.

\emph{The rough object is the volatility, not the price.} Applied to a price series
\eqref{eq:arithmeticrough} the estimator is a statement about a model this chapter has repeatedly
described as untradeable, and empirically it would return $\Hurst\approx\tfrac12$, since log-prices
are close to martingales. The series to which \eqref{eq:Hratio} is properly applied is the
log-volatility $X$ of \eqref{eq:rfsv}, where the roughness of Subsection \ref{subsec:roughvol}
actually lives and where the vol-of-vol $\nu$ cancels in the ratio exactly as $\sigma$ does here.

\emph{It is a ratio of averages, not an average of ratios.} Each expectation must be estimated across
non-overlapping windows before dividing; the estimator is consistent by the law of large numbers but
biased at finite sample, as any ratio estimator is.

\emph{It is a method of moments inside the fractional family.} Equation \eqref{eq:Hratio} is derived
for \eqref{eq:arithmeticrough} and is not model-free: it identifies $\Hurst$ on the assumption that
the series is fractional Brownian over the window. That assumption, together with the sensitivity of
all high-frequency estimators of $\Hurst$ to microstructure noise, is the subject of an active
literature, and the reported $\Hurst\approx0.1$ has been argued to be in part an artefact of it. The
estimator above inherits that debate rather than settling it.
\end{remark}

\begin{proposition}[Optimal signature schedule]\label{prop:optimalschedule}
For the model of Definition \ref{def:roughexecution} and the schedule class \eqref{eq:worked}, the
objective \eqref{eq:quadformE} separates,
\begin{equation}\label{eq:separated}
  \E\bigl[\mathcal{C}\bigr]
  =\underbrace{\eta\,\ell_\emptyset^{2}\Mat+\zeta\bigl(Q_0-\ell_\emptyset\Mat\bigr)^{2}}_{\text{schedule}}
  \;+\;\underbrace{\E\bigl[\mathbb{S}^{212}_\Mat\bigr]\,\ell_2+\mathcal{K}_{\Hurst}\,\ell_2^{2}}_{\text{signal}},
\end{equation}
with
$\mathcal{K}_{\Hurst}=\sigma^{2}\Mat^{2\Hurst+1}\bigl(\tfrac{\eta}{2\Hurst+1}+\tfrac{\zeta\Mat}{2\Hurst+2}\bigr)>0$,
and the minimisers are
\begin{equation}\label{eq:optimalpair}
  \ell_\emptyset^{*}=\frac{\zeta\,Q_0}{\eta+\zeta\Mat},
  \qquad
  \ell_2^{*}=\frac{1-2\Hurst}{4\,\eta^{\mathrm{eff}}_{\Hurst}},
  \qquad
  \eta^{\mathrm{eff}}_{\Hurst}:=\eta+\frac{(2\Hurst+1)\,\zeta\Mat}{2\Hurst+2}.
\end{equation}
The gain from the path-dependent term over the best deterministic schedule is
\begin{equation}\label{eq:gain}
  \E\bigl[\mathcal{C}(\ell^{*})\bigr]-\E\bigl[\mathcal{C}(\ell_\emptyset^{*},0)\bigr]
  \;=\;-\,\frac{(1-2\Hurst)^{2}\,\sigma^{2}\,\Mat^{2\Hurst+1}}{16\,(2\Hurst+1)\,\eta^{\mathrm{eff}}_{\Hurst}}\;\leq\;0,
\end{equation}
with equality if and only if $\Hurst=\tfrac12$.
\end{proposition}

\textbf{\textit{Proof:}} Expand the three terms of \eqref{eq:shortfallreduced} for the schedule
\eqref{eq:worked}. The inventory is $Q_t=Q_0-\ell_\emptyset t-\ell_2\mathbb{S}^{21}_t$, so
$\int_0^{\Mat}Q_t\dd\Spot_t=Q_0\mathbb{S}^{2}_\Mat-\ell_\emptyset\mathbb{S}^{12}_\Mat-\ell_2\mathbb{S}^{212}_\Mat$. Now
$\E[\mathbb{S}^{2}_\Mat]=0$, and $\E[\mathbb{S}^{12}_\Mat]=\sigma\,\E\bigl[\Mat\BH_\Mat-G_\Mat\bigr]=0$
by parts, so the first term of \eqref{eq:shortfallreduced} contributes
$\E[\mathbb{S}^{212}_\Mat]\,\ell_2$ and nothing else. The impact term is
$\eta\,\E\int_0^{\Mat}(\ell_\emptyset+\ell_2\sigma\BH_t)^{2}\dd t
=\eta\bigl(\ell_\emptyset^{2}\Mat+\ell_2^{2}\sigma^{2}\Mat^{2\Hurst+1}/(2\Hurst+1)\bigr)$, the cross
term vanishing with $\E[\BH_t]$. The terminal term is
$\zeta\,\E\bigl[(Q_0-\ell_\emptyset\Mat-\ell_2\mathbb{S}^{21}_\Mat)^{2}\bigr]
=\zeta\bigl((Q_0-\ell_\emptyset\Mat)^{2}+\ell_2^{2}\sigma^{2}\Mat^{2\Hurst+2}/(2\Hurst+2)\bigr)$,
using $\E[\mathbb{S}^{21}_\Mat]=0$ and the first identity of \eqref{eq:levelthreeaux}. Collecting gives
\eqref{eq:separated}, in which $\ell_\emptyset$ and $\ell_2$ appear in separate blocks. Each block is
a strictly convex scalar quadratic. Setting derivatives to zero,
$2\eta\Mat\ell_\emptyset=2\zeta\Mat(Q_0-\ell_\emptyset\Mat)$ gives $\ell_\emptyset^{*}$, and
$\ell_2^{*}=-\E[\mathbb{S}^{212}_\Mat]/(2\mathcal{K}_{\Hurst})$ gives, after substituting
\eqref{eq:levelthree} and simplifying, the stated form. The gain is
$-\E[\mathbb{S}^{212}_\Mat]^{2}/(4\mathcal{K}_{\Hurst})$, which is \eqref{eq:gain}. \qedfilledtext

\medskip

Three readings of Proposition \ref{prop:optimalschedule}, in increasing order of interest.

\emph{The classical answer is the $\Hurst=\tfrac12$ case.} As $\zeta\to\infty$ the terminal penalty
becomes a hard liquidation constraint and $\ell_\emptyset^{*}\to Q_0/\Mat$: the schedule is the
time-weighted average price, the standard execution benchmark and the Almgren--Chriss solution in the
absence of risk aversion \cite{almgrenchriss2001}. Simultaneously $\ell_2^{*}\to0$ when
$\Hurst=\tfrac12$. The rough machinery does not overturn the classical answer; it contains it, in the
same way that Proposition \ref{prop:purelyroughsolution} contained Black--Scholes.

\emph{The value of the signal is proportional to $(1-2\Hurst)^{2}$.} By \eqref{eq:gain} the
path-dependent term is worth exactly nothing in a Brownian market and strictly something in a rough
one, and what it is worth grows as $\Hurst$ falls away from $\tfrac12$. This is a quantitative
statement about a tradeable quantity, derived in closed form, and it is false in every model with
independent increments. It also has the right sign: for $\Hurst<\tfrac12$ we get $\ell_2^{*}>0$, that
is, sell faster when the price has risen --- which is the correct response to the mean reversion that
anti-persistence $\Hurst<\tfrac12$ encodes, and which Proposition \ref{prop:fbmproperties} b)
describes.

\emph{And the gain is not a scheduling improvement.} It is tempting to observe that no pricing
measure was used anywhere and to conclude that the obstruction of Subsection
\ref{subsec:roughexamplesolution} is irrelevant here. It is not irrelevant: it is the \emph{source}
of the gain, and the accounting is exact. Were $\Spot$ a martingale under $\PP$, we would have
$\E\int_0^{\Mat}Q_t\,\dd\Spot_t=0$ for every adapted $Q$, the first term of
\eqref{eq:shortfallreduced} would vanish identically, and
\[
  \E\bigl[\mathcal{C}\bigr]\;=\;\eta\,\E\!\int_0^{\Mat}\pi_t^{2}\,\dd t\;+\;\zeta\,\E\bigl[Q_\Mat^{2}\bigr]
\]
would not involve the price path at all. The optimal schedule would then be \emph{deterministic},
and no path-dependent rule could beat it in expectation --- which is precisely why the classical
treatment \cite{almgrenchriss2001} must adjoin a variance penalty before the problem becomes
interesting. What \eqref{eq:gain} measures is the failure of exactly that martingale property, since
$\E[\mathbb{S}^{212}_\Mat]\neq0$ \emph{is} the statement that $\E\int Q\,\dd\Spot\neq0$ for
path-dependent $Q$.

\begin{remark}[Is this an arbitrage?]\label{rem:isitarbitrage}
The model admits one, and it should be said plainly. Theorem \ref{thm:fbmnotsemimartingale} applies
to \eqref{eq:arithmeticrough} verbatim, so there is no equivalent local martingale measure and, by
Theorem \ref{thm:firstFTAP}, NFLVR fails; Rogers \cite{rogers1997} exhibits the arbitrage. Our own
computation displays a symptom of it: setting $Q_0=0$ and $\eta=\zeta=0$ in \eqref{eq:worked} leaves
the position $Q_t=-\ell_2\sigma\int_0^t\BH_s\,\dd s$, which starts from nothing and has expected
profit $-\ell_2\,\E[\mathbb{S}^{212}_\Mat]>0$ for every $\ell_2>0$ when $\Hurst<\tfrac12$, unbounded
under scaling of $\ell_2$. A positive expected return is not by itself an arbitrage, but here it is
produced by the same anti-persistence that Rogers' construction exploits.

Three consequences, which are worth keeping apart.

\emph{The execution problem remains well posed; the pricing problem would not.} This is the real
distinction, and it is sharper than the absence of a measure. Absence of arbitrage is a
\emph{hypothesis} of pricing: without an equivalent local martingale measure the valuation framework
of Subsection \ref{subsec:derivativepricing} has nothing to compute with, and the problem has no answer
rather than a wrong one. Execution is a strictly convex minimisation, and by Corollary
\ref{cor:quadformexpectation} it has a unique solution whenever $\eta,\zeta>0$, whether or not the market it is
posed in happens to be arbitrageable. For execution, absence of arbitrage is a modelling preference.

\emph{The result should be read structurally.} What Proposition \ref{prop:optimalschedule}
establishes is not a trading recommendation but a statement about separation: in a semimartingale
market, execution and alpha decouple, since expected cost is schedule-independent and scheduling is
purely a trade-off of impact against risk. Roughness couples them, and \eqref{eq:gain} is the
strength of the coupling.

\emph{The magnitude should not be trusted.} The schedule \eqref{eq:worked} rebalances continuously,
and by Cheridito \cite{cheridito2003} the arbitrage is removed by any minimum delay between
transactions, exactly as recorded in Subsection \ref{subsec:roughexamplesolution}. The economic
content survives discretisation --- an anti-persistent price mean-reverts, and that is worth
something to a trader who must sell in any case --- but the number \eqref{eq:gain} is a
continuous-time idealisation and inherits the artefact along with the closed form.
\end{remark}

\begin{remark}[Where the Hurst parameter lives, continued]\label{rem:Hlevelthree}
Remark \ref{rem:Hlivesatleveltwo} observed that $\Hurst$ first appears at level \emph{two} of the
expected signature, in the entry $\tfrac12\sigma^{2}\Mat^{2\Hurst}$, and that this is the structural
reason the estimators of Subsection \ref{subsec:roughvol} are built from second moments. Proposition
\ref{prop:levelthree} adds the next instalment, and it is a different statement. Level two knows
that the path is rough; it does not know that roughness is worth anything, because
$\tfrac12\sigma^{2}\Mat^{2\Hurst}$ is positive for every $\Hurst$ and takes no distinguished value at
$\Hurst=\tfrac12$. The first entry that vanishes exactly at $\Hurst=\tfrac12$ is the level-three
entry \eqref{eq:levelthree}, and it is precisely the coefficient of the tradeable signal in
\eqref{eq:separated}.

So the two facts sit one level apart. \emph{Roughness is a level-two phenomenon; the value of
roughness is a level-three phenomenon.} A truncation at level two would detect that the market is
rough and would find nothing to do about it. This is the concrete answer to the question of how far
up the signature one must go, and it is the reason the practical truncations quoted after Proposition
\ref{prop:factorialdecay} are $N\in\{3,4,5\}$ rather than $N=2$.
\end{remark}

\begin{remark}[What is and is not special about this example]\label{rem:executionscope}
Two honest limitations. The schedule class \eqref{eq:worked} is small --- one constant and one linear
term --- and was chosen so that \eqref{eq:optimalpair} could be written down; at larger $N$ the
solve \eqref{eq:optimalell} is still exact but is a matrix inversion rather than a formula, and the
expected-signature entries beyond level three are not available in closed form for
$\Hurst\neq\tfrac12$, as Remark \ref{rem:Hlivesatleveltwo} records. In practice $\overline{\mathbb{S}}$ is
estimated from historical paths, which is the route taken in \cite{kalsi2020}, and by Remark
\ref{rem:noliftneeded} that route needs no lift and is therefore indifferent to the restriction
$\Hurst>\tfrac13$.

What is \emph{not} special is the architecture. Lemma \ref{lem:concatenation} and Proposition
\ref{prop:shuffle} turn any objective built from time integrals, integrals against the price, and
polynomials of the control into a quadratic form in $\overline{\mathbb{S}}$; the execution problem is only
the simplest financially serious instance. The same reduction underlies signature-based pricing and
hedging of path-dependent claims \cite{lyonsnejadperez2020} and the signature market models used for
joint calibration in \cite{cuchiero2025}, which is where this circle of ideas re-enters Chapter
\ref{ch:applications}.
\end{remark}

% ---------------------------------------------------------------------------
\subsection{Two assets: the L\'evy area}\label{subsec:levyarea}

Everything so far has used a scalar asset, and by Example \ref{ex:signatureonedim} a scalar path
carries no level-two information of its own: the time coordinate had to be adjoined before the
signature said anything at all. With two assets the situation changes qualitatively, and the change
is not a matter of degree. This subsection is about the object that appears, which is the one the
whole theory is named after.

Let $X=(X^{1},X^{2})$ be the two-dimensional path of a pair of assets --- in practice
$X^{i}_{t}=\log \Spot^{i}_{t}-\log\Spot^{i}_{0}$, the cumulative log-return since the window opened
--- lifted to a geometric rough path $\rpX=(X,\lev)$. The time coordinate may be adjoined as in
Definition \ref{def:roughexecution} whenever it is wanted; it plays no role in what follows and is
suppressed. Throughout, $\theta$ denotes a lead--lag coefficient; the Heston parameter of Definition
\ref{def:roughheston} does not occur in this subsection.

\subsubsection*{One number per pair}

\begin{proposition}[The free data at level two]\label{prop:leveltwofree}
Let $\rpX$ be weakly geometric and $d$-dimensional. Then $\lev_{s,t}$ is determined by the increment
$X_{s,t}$ together with the $\tfrac12 d(d-1)$ entries $A^{ij}_{s,t}$, $i<j$, of the L\'evy area
\eqref{eq:levyarea}. In particular, for $d=1$ there are none, and for $d=2$ there is exactly one.
\end{proposition}

\textbf{\textit{Proof:}} Decompose $\lev=\Sym\lev+\Anti\lev$, the second summand being $A_{s,t}$ by
\eqref{eq:levyarea}. Proposition \ref{prop:symmetricpart} identified $\Sym\lev_{s,t}$ with
$\tfrac12X_{s,t}\otimes X_{s,t}$ for smooth paths, and \eqref{eq:weaklygeometric} is exactly the
requirement that this survive; so the symmetric part is a function of the increment alone. The
antisymmetric part is determined by its entries above the diagonal, of which there are
$\tfrac12d(d-1)$. \qedfilledtext

\medskip

Subsection \ref{subsec:tensorspaces} recorded this as the reason the second level is genuinely new
data, and Remark \ref{rem:levyarealaw} confirmed it for Brownian motion by exhibiting a law that is
not Gaussian. Read as a statement about \emph{financial} data it says something sharper. Over a
single window $\Sym\lev^{12}_{0,\Mat}=\tfrac12X^{1}_{0,\Mat}X^{2}_{0,\Mat}$ is the product of the two
total returns: it is known to anyone who has looked only at the endpoints, and averaged across
windows it is the realised covariance. Everything else level two contains --- the whole of what
depends on the \emph{trajectory} rather than on where it began and ended --- is one number, the
scalar
\begin{equation}\label{eq:pairarea}
  A_{s,t}\;=\;\tfrac12\Bigl(\lev^{12}_{s,t}-\lev^{21}_{s,t}\Bigr)
  \;=\;\tfrac12\left(\int_s^t X^{1}_{s,r}\,\dd X^{2}_{r}-\int_s^t X^{2}_{s,r}\,\dd X^{1}_{r}\right)
\end{equation}
of Remark \ref{rem:levyarealaw}, now carried by a pair of assets rather than by a planar Brownian
motion. It vanishes when the two components move proportionally, a straight line enclosing no area.

\begin{remark}[The area is a profit and loss]\label{rem:areaispnl}
The reading that matters here is not geometric. The integral $\int_s^tX^{1}_{s,r}\,\dd X^{2}_{r}$ is
the profit of a strategy holding, at each instant, a position in asset $2$ equal to the move asset
$1$ has made since the window opened; and symmetrically for the other term. So $A_{s,t}$
is one half the difference of two crossed pairs strategies --- long the laggard against the leader,
in whichever direction the pair has been moving --- and it is a quantity a desk can actually realise.
Level-two information is not an abstraction: for a pair of assets it \emph{is} the return of a
specific long--short book.
\end{remark}

\begin{remark}[Two conventions that do not matter, and one that does]\label{rem:areaconventions}
The area is insensitive to the It\^o--Stratonovich choice. Passing from one to the other adds
$\tfrac12\qvv{X^{1}}{X^{2}}$ to each of the two integrals in \eqref{eq:levyarea}, and the two
additions cancel in the difference; the ambiguity of Chapter \ref{ch:foundations} lives entirely in
$\Sym\lev$, which is the quadratic covariation, and not in $\Anti\lev$. What the area \emph{is}
sensitive to is the choice of lift itself: by Proposition \ref{prop:liftnotunique} a path does not
determine its own second level, and by Proposition \ref{prop:leveltwofree} the whole of that freedom,
for a pair, is the single number $A$. Two markets with identical realised price paths and
identical realised covariances but different areas are different rough paths, and by the Universal
Limit Theorem \ref{thm:ult} they produce different solutions to any rough differential equation
driven by them --- that is, different profit and loss for any path-dependent book. Correlation is not
merely an incomplete description of a pair; it is provably the part of level two that carries nothing
new.
\end{remark}

\subsubsection*{Aggregation is not additive}

Returns add across adjacent windows. Areas do not, and the defect is exactly Chen's relation.

\begin{proposition}[Chen's relation for the area]\label{prop:chenarea}
For $s\leq u\leq t$,
\begin{equation}\label{eq:chenarea}
  A_{s,t}\;=\;A_{s,u}+A_{u,t}
  \;+\;\tfrac12\det\begin{pmatrix} X^{1}_{s,u} & X^{1}_{u,t}\\[2pt] X^{2}_{s,u} & X^{2}_{u,t}\end{pmatrix}.
\end{equation}
\end{proposition}

\textbf{\textit{Proof:}} Chen's relation \eqref{eq:chenaxiom} reads
$\lev_{s,t}=\lev_{s,u}+\lev_{u,t}+X_{s,u}\otimes X_{u,t}$. Taking the antisymmetric part of the
$(1,2)$ entry, the last term contributes
$\tfrac12\bigl(X^{1}_{s,u}X^{2}_{u,t}-X^{2}_{s,u}X^{1}_{u,t}\bigr)$, which is the stated
determinant. \qedfilledtext

\medskip

The correction term is the signed area of the triangle spanned by the two chord returns, and it is
zero exactly when the two sub-period returns are collinear. Financially: the profit of the pairs book
of Remark \ref{rem:areaispnl} over a month is \emph{not} the sum of its profits over the constituent
weeks, and the discrepancy is a determinant of weekly returns. Nothing in a return-based description
of a pair has an analogue of \eqref{eq:chenarea}; it is a structural constraint on how path
information aggregates, and it is exact rather than approximate.

\subsubsection*{A computable lead--lag model}

\begin{definition}[Lead--lag pair]\label{def:leadlag}
Let $\BH$ and $\widetilde{W}^{\Hurst}$ be independent scalar fractional Brownian motions with the same
$\Hurst\in(\tfrac13,\tfrac12]$, let $\sigma_{1},\sigma_{2}>0$, $\rho\in[-1,1]$ and $\theta\in\R$, and
set
\begin{equation}\label{eq:leadlagmodel}
  X^{1}_{t}=\sigma_{1}\BH_{t},
  \qquad
  X^{2}_{t}=\sigma_{2}\Bigl(\rho\,\BH_{t}+\sqrt{1-\rho^{2}}\,\widetilde{W}^{\Hurst}_{t}\Bigr)
  \;+\;\theta\int_{0}^{t}\BH_{s}\,\dd s .
\end{equation}
The second asset moves with the first contemporaneously, through $\rho$, and responds to its
\emph{past} through $\theta$: asset $1$ leads.
\end{definition}

\begin{proposition}[Expected area of a lead--lag pair]\label{prop:leadlagarea}
For the pair \eqref{eq:leadlagmodel},
\begin{equation}\label{eq:leadlagarea}
  \E\bigl[A_{0,\Mat}\bigr]
  \;=\;\frac{\sigma_{1}\,\theta\,(3-2\Hurst)\,\Mat^{2\Hurst+1}}{4\,(2\Hurst+1)} .
\end{equation}
Neither $\rho$ nor $\sigma_{2}$ appears.
\end{proposition}

\textbf{\textit{Proof:}} Write $G_{t}=\int_{0}^{t}\BH_{s}\,\dd s$ and expand the two integrals of
\eqref{eq:levyarea}. Using $\dd X^{2}=\sigma_{2}\rho\,\dd\BH+\sigma_{2}\sqrt{1-\rho^{2}}\,\dd\widetilde
W^{\Hurst}+\theta\,\BH_{t}\,\dd t$,
\[
  \int_{0}^{\Mat}X^{1}\,\dd X^{2}
  =\sigma_{1}\sigma_{2}\rho\!\int_{0}^{\Mat}\!\BH\dd\BH
  +\sigma_{1}\sigma_{2}\sqrt{1-\rho^{2}}\!\int_{0}^{\Mat}\!\BH\,\dd\widetilde W^{\Hurst}
  +\sigma_{1}\theta\!\int_{0}^{\Mat}\!\bigl(\BH_{t}\bigr)^{2}\dd t,
\]
while $\dd X^{1}=\sigma_{1}\dd\BH$ gives
\[
  \int_{0}^{\Mat}X^{2}\,\dd X^{1}
  =\sigma_{1}\sigma_{2}\rho\!\int_{0}^{\Mat}\!\BH\dd\BH
  +\sigma_{1}\sigma_{2}\sqrt{1-\rho^{2}}\!\int_{0}^{\Mat}\!\widetilde W^{\Hurst}\dd\BH
  +\sigma_{1}\theta\!\int_{0}^{\Mat}\!G_{t}\,\dd\BH_{t}.
\]
The two $\rho$ terms are \emph{identical} and cancel in the difference \eqref{eq:levyarea}. The two
mixed terms have expectation zero: the law of $\widetilde W^{\Hurst}$ is invariant under
$\widetilde W^{\Hurst}\mapsto-\widetilde W^{\Hurst}$, which commutes with the piecewise-linear
approximations defining the joint geometric lift and reverses the sign of each mixed integral, so
each is symmetrically distributed. There remain
$\E\int_{0}^{\Mat}(\BH_{t})^{2}\dd t=\Mat^{2\Hurst+1}/(2\Hurst+1)$ and, by Proposition
\ref{prop:levelthree} with $\sigma=1$,
$\E\int_{0}^{\Mat}G_{t}\dd\BH_{t}=(2\Hurst-1)\Mat^{2\Hurst+1}/\bigl(2(2\Hurst+1)\bigr)$. Hence
\[
  \E\bigl[A_{0,\Mat}\bigr]
  =\frac{\sigma_{1}\theta}{2}\,\Mat^{2\Hurst+1}\left(\frac{1}{2\Hurst+1}-\frac{2\Hurst-1}{2(2\Hurst+1)}\right)
  =\frac{\sigma_{1}\theta\,\Mat^{2\Hurst+1}}{2(2\Hurst+1)}\cdot\frac{3-2\Hurst}{2},
\]
which is \eqref{eq:leadlagarea}. \qedfilledtext

\medskip

At $\Hurst=\tfrac12$ the expression reduces to $\sigma_{1}\theta\Mat^{2}/4$, which is checked
directly: there $\int G\,\dd\Wt$ is an It\^o integral of mean zero and only the first term survives.
At $\theta=0$ it vanishes for every $\rho$, as it must, the L\'evy area of a correlated Brownian pair
having mean zero.

\begin{remark}[Correlation is invisible to the area]\label{rem:rhocancels}
The cancellation of $\rho$ in the proof is the whole content of this subsection, written as a
formula. Proposition \ref{prop:leveltwofree} said that the symmetric part of level two is the
endpoints and the antisymmetric part is everything else; Proposition \ref{prop:leadlagarea} exhibits
a pair in which the contemporaneous coupling $\rho$ and the lagged coupling $\theta$ are separately
adjustable, and shows that the area sees the second and is \emph{exactly blind} to the first. A
correlation estimate and an area estimate are therefore not two views of the same quantity; they are
measurements of orthogonal components of the same tensor.

Inverting \eqref{eq:leadlagarea} estimates the lead--lag coefficient,
\[
  \widehat\theta\;=\;\frac{4(2\Hurst+1)\,\widehat{\E}\bigl[A_{0,\Mat}\bigr]}
  {\sigma_{1}(3-2\Hurst)\,\Mat^{2\Hurst+1}},
\]
from a single level-two entry. The classical route to the same conclusion is a cross-correlogram
scanned over a grid of lags; \eqref{eq:levyarea} aggregates that scan into one number, at the cost of
committing to the shape \eqref{eq:leadlagmodel} in exchange for not having to choose a maximum lag.
\end{remark}

\begin{remark}[What is computable, and how]\label{rem:areacomputable}
Every object in this subsection is a Riemann--Stieltjes sum on sampled data. With observations at
$0=t_{0}<\dots<t_{n}=\Mat$,
\[
  \lev^{12}_{0,\Mat}\;\approx\;\sum_{k}X^{1}_{0,t_{k}}\bigl(X^{2}_{t_{k+1}}-X^{2}_{t_{k}}\bigr),
\]
and $A$ is half the difference of that and its transpose. By Remark
\ref{rem:noliftneeded} no lift is involved: the interpolated data has bounded variation. Two exact
identities are available as checks on any implementation, and both hold to machine precision rather
than statistically --- antisymmetry, $A^{12}=-A^{21}$, and geometricity,
$\lev^{12}_{0,\Mat}+\lev^{21}_{0,\Mat}=X^{1}_{0,\Mat}X^{2}_{0,\Mat}$, the latter being Proposition
\ref{prop:leveltwofree} read as an assertion about numbers. Chen's relation \eqref{eq:chenarea}
supplies a third, across any splitting of the window.
\end{remark}

\section{The Boundary of the Theory}\label{sec:beyond}

It would be a poor advertisement for a chapter on roughness to conclude without saying where its
own methods give out, particularly since Remark \ref{rem:H01outside} has already located the point.

Rough path theory lifts fractional Brownian motion for $\Hurst > \tfrac14$. The volatility data
suggest $\Hurst \approx 0.1$. The gap is not an artefact of the level-two restriction
\eqref{eq:alpharange} that we adopted for clarity: for $\Hurst \leq \tfrac14$ the piecewise-linear
approximations fail to converge at \emph{any} level, and the canonical lift does not exist.

The difficulty is visible in the pricing problem itself. Writing $\widehat W_t = \int K^{\Hurst}(t-s)\,\dd W_s$
for the Riemann--Liouville form of the driving noise in \eqref{eq:rfsv}, with kernel
$K^{\Hurst}(u) = u^{\Hurst - 1/2}\mathds{1}_{u>0}$, one must make sense of integrals of the type
\begin{equation}\label{eq:roughvolintegral}
\int_0^{\Mat} F\bigl(\widehat W_t\bigr)\,\dd W_t .
\end{equation}
The integrator $W$ is a perfectly good Brownian rough path. The integrand, however, has Hölder
regularity $\Hurst^{-}$, which for $\Hurst<\tfrac12$ is \emph{worse} than that of $W$, so
$F(\widehat W)$ is not a path controlled by $W$ in the sense of Definition \ref{def:controlled}, and
Theorem \ref{thm:gubinelli} does not apply. The Taylor expansion \eqref{eq:controlleddef} is simply
the wrong ansatz.

It is worth locating \eqref{eq:roughvolintegral} precisely, because without it the model is not hard
at all. Let the price be driven by $\Wt^{\Spot}=\rho\,\Wt+\sqrt{1-\rho^{2}}\,\Wt^{\perp}$, with
$\Wt$ the Brownian motion generating $\widehat W$ and $\Wt^{\perp}$ independent of it, and condition
on the volatility path. This is the device of Romano and Touzi \cite{romanotouzi1997}: the value of a
European claim becomes a mixture of Black--Scholes values \eqref{eq:bscall}, taken over the realised
variance, and the mixture depends on the trajectory through exactly two quantities. The first is the
realised variance $\int_0^{\Mat}v_s\,\dd s$, an ordinary Riemann integral of a continuous process,
which requires no integration theory whatever. The second is the \emph{leverage} term
\begin{equation}\label{eq:leverageterm}
  \rho\int_0^{\Mat}\sqrt{v_s}\,\dd\Wt_s,\qquad v_s=F\bigl(\widehat W_s\bigr),
\end{equation}
which is \eqref{eq:roughvolintegral} with $F$ replaced by $F^{1/2}$. At $\rho=0$ the term is absent
and a rough volatility model is analytically no harder than a model with deterministic volatility.
Every difficulty catalogued in this section lives in \eqref{eq:leverageterm} and nowhere else.

That localisation is not a technicality, because the correlation is not a technicality. A negative
$\rho$ is what produces the skew of Subsection \ref{subsec:smile_skew}, and by
\eqref{eq:skewpowerlaw} it is what produces the short-dated power law that motivated Part
\ref{part:part2} in the first place. The feature which makes a rough volatility model worth having
is precisely the feature which puts it beyond the integration theory of Chapter
\ref{ch:foundations}.

What replaces it is Hairer's theory of \emph{regularity structures}, of which rough path theory is
the one-dimensional special case. There, the fixed local expansion \eqref{eq:controlleddef} in
powers of $X_{s,t}$ is replaced by an expansion in an abstract, problem-dependent basis of symbols,
and the products which fail to converge are repaired by \emph{renormalisation} --- subtracting
divergent constants in a manner constrained by the algebra rather than chosen by hand. Applied to
\eqref{eq:roughvolintegral} the construction requires four symbols, corresponding to the powers
$(\widehat W_{s,t})^m$ for $m \leq 3$, and yields a renormalised Wong--Zakai theorem of the form
\begin{equation}\label{eq:renormalisedwz}
\int_0^t F\bigl(\widehat W^{\varepsilon}\bigr)\,\dd W^{\varepsilon} \;-\; c^{\varepsilon}\!\int_0^t F'\bigl(\widehat W^{\varepsilon}\bigr)\,\dd s
\;\longrightarrow\; \int_0^t F\bigl(\widehat W\bigr)\,\dd W ,
\end{equation}
valid for $\Hurst > \tfrac18$ and with a constant $c^\varepsilon$ which \emph{diverges} as the
mollification parameter $\varepsilon$ tends to zero.

The reader who has been keeping score will object that $\tfrac18 = 0.125$ exceeds the
$\Hurst \approx 0.1$ reported in Subsection \ref{subsec:roughvol}, so that this construction misses
the empirical exponent as well. That is correct, and the resolution is favourable: the threshold
$\tfrac18$ is an artefact of stopping at four symbols, not a limit of the method. The number of
symbols required is governed by the smallest $M$ with $(M+1)\Hurst + \tfrac12 > 1$, so admitting
more symbols lowers the threshold, and the construction extends to every
$\Hurst \in (0,\tfrac12)$ at the cost of a regularity structure whose size grows as
$\Hurst \downarrow 0$. Roughness is therefore expensive rather than prohibitive --- which is the
honest summary, and a different statement from the genuine obstruction of Remark
\ref{rem:H01outside}, where no number of levels suffices because the lift itself fails to exist. At $\Hurst = \tfrac12$ the constant is
$\tfrac12$ and \eqref{eq:renormalisedwz} degenerates into the Itô--Stratonovich correction of
Chapter \ref{ch:foundations}, which is a pleasing consistency check and a fair summary of the whole
subject: the corrections that stochastic calculus treats as fixed constants are, at lower
regularity, divergent quantities held in place by an algebraic structure.

It is worth being exact about what this machinery is needed \emph{for}, since the honest answer is
narrower than the apparatus suggests, and sharper than a general appeal to roughness. The integrand
$F(\widehat W)$ is adapted to the filtration of $\Wt$ and square-integrable, so
\eqref{eq:roughvolintegral} exists as an ordinary It\^o integral by Section
\ref{sec:stochasticcalculus} and nothing further: neither rough paths nor regularity structures are
required in order to \emph{define} a rough volatility model, nor to simulate one. What they are
required for is threefold, and each item is load-bearing.

\begin{itemize}
  \item[i)] \emph{A pathwise construction.} The It\^o integral is an $L^{2}$ limit defined off a null
    set, and is not a continuous function of the trajectory of $\Wt$; no statement about an
    individual realisation is available from it.
  \item[ii)] \emph{The limit of a discretisation.} Equation \eqref{eq:renormalisedwz} says that a
    smoothed approximation converges to the wrong object unless the divergent counterterm is
    subtracted. A naive scheme can therefore converge, cleanly and without any numerical symptom, to
    a model other than the one intended. This is the item with practical consequences.
  \item[iii)] \emph{Asymptotics.} The contraction principle requires a continuous map, so a
    small-time large deviation principle for the price --- and with it the short-dated smile and the
    exponent of \eqref{eq:skewpowerlaw} --- becomes available only after the lift. It is here, rather
    than in the construction of the model, that the continuity of the It\^o--Lyons map, Theorem
    \ref{thm:ult}, does financial work, and it is the answer to the debt acknowledged in Remark
    \ref{rem:roughhestoncost}.
\end{itemize}

We shall not develop any of this. The regularity structure for rough volatility is constructed in
\cite{bayerfrizgassiat2020}, from which the count of symbols and the threshold $\Hurst>\tfrac18$
quoted above are taken; the general theory is set out in \cite[Ch.\ 13--14]{frizhairer2020}, and the
rough volatility literature built on it is surveyed in \cite{bayerfrizfukasawa2023} and
\cite{dinunno2023}. It is recorded here to mark honestly the edge of what has been proved. The applications of Chapter
\ref{ch:applications} are constructed so as not to depend on a lifting theorem that is unavailable at the
exponent the data demand --- the local volatility surface of Subsection \ref{subsec:localvol} being
a classical diffusion, and the rough Heston model of Definition \ref{def:roughheston} being
tractable through the affine route of Theorem \ref{thm:roughhestoncf} rather than through rough
integration.

% Unnumbered, so deliberately unlabelled: a \label here would bind to whatever
% counter was last stepped and any \ref to it would print the wrong number.
\section*{Coda: the three re-derivations}
\addcontentsline{toc}{section}{Coda: the three re-derivations}

It is worth gathering, before closing, the promise made at the end of Section
\ref{sec:roughintro}. Three results of Chapter \ref{ch:foundations} have been obtained again in this
chapter on wholly different foundations, and the accounting in each case was stated where the result
appeared:

\begin{itemize}
    \item[i)] \textbf{Itô's formula} (Corollary \ref{cor:itoformula}). Recovered from the rough
    chain rule \eqref{eq:roughchainrule} applied to the Itô lift, whose bracket
    \eqref{eq:bracket} is $t\,\mathrm{Id}$. \emph{Cost:} $F \in \mathcal{C}^3_b$ rather than
    $\theta \in \mathcal{C}^{1,2}$, and time-homogeneity. \emph{Gain:} the identity holds pathwise,
    for a fixed trajectory and its lift, with both sides continuous in $\rpdist$.
    \item[ii)] \textbf{The Itô--Stratonovich correction} \eqref{eq:integraldifference}. Recovered
    not as a theorem about Brownian motion but as the general formula for the difference between
    two lifts of the \emph{same} path, evaluated at $f_t = \tfrac12 t\,\mathrm{Id}$.
    \emph{Cost:} none of substance. \emph{Gain:} the correction is exhibited as a consequence of
    choosing $\levW^{\mathrm{It\hat o}}$ over $\levW^{\mathrm{Strat}}$, rather than as a feature of
    the integral.
    \item[iii)] \textbf{Wong--Zakai} (Corollary \ref{cor:wongzakai}). Recovered in two lines from
    the continuity of the Itô--Lyons map, Theorem \ref{thm:ult}. \emph{Cost:} the lifting theorem
    for $W$, performed once in Example \ref{ex:brownianroughpath}. \emph{Gain:} the reason the
    limit is Stratonovich and not Itô becomes visible --- smooth approximations carry geometric
    lifts, and the limit remembers the area.
\end{itemize}

Taken together these say something the classical theory cannot. In the classical account --- that of
Chapter \ref{ch:foundations}, extended by the Wong--Zakai theorem --- the second-order term of Itô's
rule, the $\tfrac12\int F'$ of the Itô--Stratonovich dictionary, and the failure of smooth
approximation to converge to the Itô solution are three separate facts,
each proved by its own argument and each attributable to the non-vanishing quadratic variation of
Brownian motion. Here they are one fact, seen three times: the trajectory does not determine its own
iterated integrals, a choice must be made, and every one of the three phenomena is that choice
becoming visible. Equation \eqref{eq:renormalisedwz} above is the same statement pushed to lower
regularity, where the constant held in place by the algebra is no longer $\tfrac12$ but divergent.
